% Chapter 3, Section 1 _Linear Algebra_ Jim Hefferon
%  http://joshua.smcvt.edu/linearalgebra
%  2001-Jun-11
\chapter{Pemetaan Antar-Ruang}% 
% Without this hack, it prints \nobreak
\ \vspace{-10ex}

\index{isomorfisme|(}
\section{Isomorfisme} %
Dalam contoh-contoh setelah definisi ruang vektor,
kita mengemukakan intuisi bahwa beberapa ruang pada hakikatnya sama dengan
ruang lainnya.
Sebagai contoh, 
% the space of two-tall column vectors and
% the space of two-wide row vectors
% are not equal because their 
% elements\Dash column vectors and row vectors\Dash are 
% not equal, 
kita dapat melihat bahwa ruang matriks dua-kali-dua dan
ruang vektor kolom berukuran empat
\begin{equation*}
  \matspace_{\nbyn{2}}=\set{
    \begin{pmatrix}
      a &b \\
      c &d
    \end{pmatrix}\suchthat a,b,c,d\in\R}
  \qquad
  \R^4=\set{\colvec{a \\ b \\ c \\ d}\suchthat a,b,c,d\in\R}
\end{equation*}
tidak sama karena unsur-unsurnya tidak sama, tetapi perbedaannya
hanyalah pada tampilan.
Sekarang kita akan menyatakan hal ini secara tepat.

Hal ini menggambarkan tahap yang umum dalam penyelidikan matematis.
Dengan bantuan beberapa contoh, kita telah memperoleh suatu gagasan.
Selanjutnya kita akan memberikan definisi formal, lalu
menghasilkan sejumlah hasil yang mendukung pendapat kita bahwa
definisi tersebut menangkap gagasan itu.
Kita sudah pernah melihat pola ini, misalnya pada bagian pertama
dalam bab Ruang Vektor.
Di sana, kajian tentang sistem linear mengarahkan kita untuk
meninjau kumpulan yang tertutup terhadap kombinasi linear.
Kumpulan semacam itu kita definisikan sebagai ruang vektor,
lalu kita memberikan beberapa hasil pendukung.

Itu bukanlah titik akhir; sebaliknya,
hal tersebut mengarah pada wawasan baru seperti gagasan tentang basis.
Di sini pun, setelah membuat dan mendukung sebuah definisi,
kita akan memperoleh dua kejutan yang menyenangkan.
Pertama, kita akan mendapati bahwa definisi tersebut berlaku pada
beberapa kasus menarik yang tidak terduga.
Kedua, kajian terhadap definisi itu akan menghasilkan gagasan-gagasan baru.
Dengan demikian, penyelidikan kita akan terus berkembang.








\subsection{Def{}inisi dan Contoh}
Kita mulai dengan dua contoh yang mengarahkan kita pada definisi yang tepat. 

\begin{example}  \label{exam:TwoWideIsoTwoTall}
Ruang vektor baris berukuran dua dan
ruang vektor kolom berukuran dua
``sama'' dalam arti bahwa jika kita memasangkan vektor-vektor
yang memiliki komponen yang sama, misalnya,
\begin{equation*}
   \rowvec{1 &2} \quad\longleftrightarrow\quad \colvec[r]{1 \\ 2}
\end{equation*}
(baca panah ganda sebagai ``berkorespondensi dengan''),
maka pemasangan ini menghormati operasi-operasinya.
Sebagai contoh, vektor-vektor yang berkorespondensi ini dijumlahkan menjadi
hasil-hasil yang juga berkorespondensi
\begin{equation*}
  \rowvec{1 &2}+\rowvec{3 &4}=\rowvec{4 &6}
  \quad\longleftrightarrow\quad
  \colvec[r]{1 \\ 2}+\colvec[r]{3 \\ 4}=\colvec[r]{4 \\ 6}
\end{equation*}
dan berikut adalah contoh bahwa korespondensi tersebut menghormati
perkalian skalar.
\begin{equation*}
  5\cdot\rowvec{1  &2}=\rowvec{5 &10}
  \quad\longleftrightarrow\quad
  5\cdot\colvec[r]{1 \\ 2}=\colvec[r]{5 \\ 10}
\end{equation*}
Secara umum, di bawah korespondensi
\begin{equation*}
  \rowvec{a_0  &a_1}
  \quad\longleftrightarrow\quad
  \colvec{a_0 \\ a_1}
\end{equation*} 
kedua operasi dipertahankan:
\begin{equation*}
  \rowvec{a_0  &a_1}+\rowvec{b_0  &b_1}=\rowvec{a_0+b_0 &a_1+b_1}
  \hfill\longleftrightarrow\hfill
  \colvec{a_0 \\ a_1}+\colvec{b_0  \\ b_1}=\colvec{a_0+b_0  \\ a_1+b_1}
\end{equation*}
dan
\begin{equation*}
 r\cdot\rowvec{a_0 &a_1}=\rowvec{ra_0 &ra_1}
  \quad\longleftrightarrow\quad
 r\cdot\colvec{a_0  \\ a_1}=\colvec{ra_0  \\ ra_1}
\end{equation*}
(semua variabelnya adalah skalar).
\end{example}

\begin{example} \label{exam:PolyTwoIsoRThree}  
Dua ruang lain yang dapat kita pandang ``sama'' adalah
\( \polyspace_2 \), ruang polinomial kuadrat, dan \( \Re^3 \).
Korespondensi yang alami adalah sebagai berikut.
\begin{equation*}
  a_0+a_1x+a_2x^2
  \quad\longleftrightarrow\quad
  \colvec{a_0 \\ a_1 \\ a_2}
%  \qquad
  \tag*{\mbox{(mis., $1+2x+3x^2\,\longleftrightarrow\,\colvec[r]{1 \\ 2 \\ 3}$)}}
\end{equation*}
Korespondensi ini mempertahankan struktur:~unsur-unsur yang berkorespondensi
dijumlahkan dengan cara yang berkorespondensi
\begin{equation*}
    \mbox{\begin{tabular}{r}
        \( a_0+a_1x+a_2x^2 \)  \\
     +\,\,\( b_0+b_1x+b_2x^2 \)  \\ \hline
        \( (a_0+b_0)+(a_1+b_1)x+(a_2+b_2)x^2 \)
     \end{tabular} }
    \longleftrightarrow
    \colvec{a_0 \\ a_1 \\ a_2}
    +\colvec{b_0 \\ b_1 \\ b_2}
    =\colvec{a_0+b_0 \\ a_1+b_1 \\ a_2+b_2}
\end{equation*}
dan perkalian skalarnya juga berkorespondensi.
\begin{equation*}
  r\cdot(a_0+a_1x+a_2x^2)=
    (ra_0)+(ra_1)x+(ra_2)x^2
  \quad\longleftrightarrow\quad
  r\cdot\colvec{a_0 \\ a_1 \\ a_2}
  =\colvec{ra_0 \\ ra_1 \\ ra_2}
\end{equation*}
\end{example}

\begin{definition} \label{def:Isomorphism}
%<*df:Isomorphism>
Sebuah \definend{isomorfisme}\index{isomorfisme!definisi}%
\index{ruang vektor!isomorfisme}
antara dua ruang vektor $V$ dan $W$
adalah pemetaan \( \map{f}{V}{W} \) yang
\begin{enumerate}
  \item merupakan korespondensi:\index{korespondensi}
     \( f \) bersifat injektif dan surjektif;% \footnotemark % keep fn in minipage
                   \appendrefs{korespondensi}
  \item \definend{mempertahankan struktur:}\index{fungsi!mempertahankan struktur}
   jika \( \vec{v}_1,\vec{v}_2\in V \), maka
   \begin{equation*}
      f(\vec{v}_1+\vec{v}_2)=f(\vec{v}_1)+f(\vec{v}_2)
   \end{equation*}
   dan jika \( \vec{v}\in V \) dan \( r\in\Re \), maka
   \begin{equation*}
         f(r\vec{v})=r{}f(\vec{v})  % suppress kerning
   \end{equation*}
\end{enumerate}
(kita menulis \( V\isomorphicto W \), dibaca ``\( V \) isomorfik dengan \( W \)'',
jika pemetaan semacam itu ada).
%</df:Isomorphism>
\end{definition}% \footnotetext{More information on one-to-one and onto maps 
                %               is in the appendix.}

\noindent ``Morfisme''\index{morfisme} berarti pemetaan,
sehingga ``isomorfisme'' berarti pemetaan yang menyatakan kesamaan.

\begin{example}  \label{ex:CheckMapIsIso}
Ruang vektor
\(  G=\set{c_1\cos\theta+c_2\sin\theta\suchthat c_1,c_2\in\Re} \)
yang terdiri atas fungsi-fungsi dari \( \theta \) isomorfik dengan
\( \Re^2 \) di bawah pemetaan berikut.
\begin{equation*}
  c_1\cos\theta+c_2\sin\theta\mapsunder{f}\colvec{c_1 \\ c_2}
\end{equation*}
Kita akan memeriksanya dengan menelusuri syarat-syarat dalam definisi.
Pertama-tama kita akan memverifikasi syarat~(1), yaitu bahwa pemetaan tersebut
merupakan korespondensi antara himpunan-himpunan yang mendasari kedua ruang.

Untuk menunjukkan bahwa $f$ injektif,
kita harus membuktikan bahwa \( f(\vec{a})=f(\vec{b}) \) hanya jika
\( \vec{a}=\vec{b} \).
Jika
\begin{equation*}
   f(a_1\cos\theta+a_2\sin\theta)=f(b_1\cos\theta+b_2\sin\theta)
\end{equation*}
maka berdasarkan definisi $f$
\begin{equation*}
  \colvec{a_1 \\ a_2}=\colvec{b_1 \\ b_2}
\end{equation*}
dari sini kita menyimpulkan bahwa \( a_1=b_1 \) dan \( a_2=b_2 \),
karena dua vektor kolom sama hanya jika komponen-komponennya sama.
Jadi $a_1\cos\theta+a_2\sin\theta=b_1\cos\theta+b_2\sin\theta$, dan
sesuai kebutuhan kita telah memverifikasi bahwa \( f(\vec{a})=f(\vec{b}) \)
menyiratkan \( \vec{a}=\vec{b} \).
% which shows that \( f \) is one-to-one.

Untuk membuktikan bahwa $f$ surjektif, kita harus memeriksa bahwa setiap unsur
kodomain \( \Re^2 \) merupakan bayangan dari suatu unsur domain $G$.
Maka, tinjau suatu unsur kodomain
\begin{equation*}
   \colvec{x \\ y}
\end{equation*}
dan perhatikan bahwa unsur ini merupakan bayangan
\( x\cos\theta+y\sin\theta \) di bawah $f$.

Selanjutnya kita akan memverifikasi syarat~(2), yaitu bahwa $f$
mempertahankan struktur.
Perhitungan berikut menunjukkan bahwa \( f \) mempertahankan penjumlahan.
\begin{multline*}
  f\bigl(\,(a_1\cos\theta+a_2\sin\theta)
      +(b_1\cos\theta+b_2\sin\theta)\,\bigr)                     \\
 \begin{aligned}
  &=f\bigl(\,(a_1+b_1)\cos\theta+(a_2+b_2)\sin\theta\,\bigr)   \\
  &=\colvec{a_1+b_1 \\ a_2+b_2}   \\
  &=\colvec{a_1 \\ a_2}+\colvec{b_1 \\ b_2}   \\
  &=f(a_1\cos\theta+a_2\sin\theta)+f(b_1\cos\theta+b_2\sin\theta)
 \end{aligned}
\end{multline*}
Perhitungan yang menunjukkan bahwa $f$ mempertahankan perkalian skalar serupa.
\begin{align*}
  f\bigl(\,r\cdot(a_1\cos\theta+a_2\sin\theta)\,\bigr)
  &=f(\,ra_1\cos\theta+ra_2\sin\theta\,)   \\
  &=\colvec{ra_1 \\ ra_2}   \\
  &=r\cdot\colvec{a_1 \\ a_2}   \\
  &=r\cdot\, f(a_1\cos\theta+a_2\sin\theta)
\end{align*}

Karena (1) dan~(2) telah diverifikasi,
kita mengetahui bahwa $f$ adalah isomorfisme dan dapat mengatakan bahwa
kedua ruang tersebut isomorfik, $G\isomorphicto\Re^2$.
\end{example}

\begin{example} \label{ex:LinComboThreeIsoPTwo}
Misalkan \( V \) adalah ruang
\( \set{c_1x+c_2y+c_3z\suchthat c_1,c_2,c_3\in\Re} \)
yang terdiri atas kombinasi linear ketiga variabel tersebut,
di bawah operasi penjumlahan dan perkalian skalar yang alami.
Maka $V$ isomorfik dengan $\polyspace_2$, yaitu ruang polinomial kuadrat.

Untuk menunjukkan hal ini, kita harus menghasilkan sebuah pemetaan isomorfisme.
Terdapat lebih dari satu kemungkinan; misalnya, berikut empat pilihan.
\begin{center}
  \begin{tabular}{c}
    $c_1x+c_2y+c_3z$
  \end{tabular}  \quad
  \begin{tabular}{rl}
      $\mapsunder{f_1}$    &$c_1+c_2x+c_3x^2$  \\
      $\mapsunder{f_2}$    &$c_2+c_3x+c_1x^2$  \\
      $\mapsunder{f_3}$    &$-c_1-c_2x-c_3x^2$  \\
      $\mapsunder{f_4}$    &$c_1+(c_1+c_2)x+(c_1+c_3)x^2$
  \end{tabular}
\end{center}
Pemetaan pertama merupakan korespondensi yang lebih alami karena hanya
memindahkan koefisien-koefisiennya.
Namun, kita akan menggunakan $f_2$ untuk menegaskan bahwa terdapat
isomorfisme selain yang tampak jelas.
(Pemeriksaan bahwa $f_1$ merupakan isomorfisme terdapat pada
\nearbyexercise{exer:NatMapAlsoIso}.)

Untuk menunjukkan bahwa $f_2$ injektif, kita akan membuktikan bahwa jika
$f_2(c_1x+c_2y+c_3z)=f_2(d_1x+d_2y+d_3z)$
maka $c_1x+c_2y+c_3z=d_1x+d_2y+d_3z$.
Asumsi $f_2(c_1x+c_2y+c_3z)=f_2(d_1x+d_2y+d_3z)$,
berdasarkan definisi $f_2$, memberikan
$c_2+c_3x+c_1x^2=d_2+d_3x+d_1x^2$.
Dua polinomial yang sama memiliki koefisien yang sama, sehingga
$c_2=d_2$, $c_3=d_3$, dan $c_1=d_1$.
Oleh karena itu, $f_2(c_1x+c_2y+c_3z)=f_2(d_1x+d_2y+d_3z)$ menyiratkan
$c_1x+c_2y+c_3z=d_1x+d_2y+d_3z$, sehingga $f_2$ injektif.

Pemetaan $f_2$ surjektif karena setiap unsur $a+bx+cx^2$ dari kodomain
merupakan bayangan suatu unsur domain, tepatnya $f_2(cx+ay+bz)$.
Sebagai contoh, $2+3x-4x^2$ adalah $f_2(-4x+2y+3z)$.

Perhitungan untuk mempertahankan struktur serupa dengan perhitungan pada
contoh sebelumnya.
Pemetaan $f_2$ mempertahankan penjumlahan
\begin{multline*}
  f_2\bigl((c_1x+c_2y+c_3z)
        +(d_1x+d_2y+d_3z)\bigr)       \\
 \begin{aligned}
  &=f_2\bigl((c_1+d_1)x+(c_2+d_2)y+(c_3+d_3)z\bigr)   \\
  &=(c_2+d_2)+(c_3+d_3)x+(c_1+d_1)x^2   \\
  &=(c_2+c_3x+c_1x^2)+(d_2+d_3x+d_1x^2)   \\
  &=f_2(c_1x+c_2y+c_3z)+f_2(d_1x+d_2y+d_3z)
 \end{aligned}
\end{multline*}
dan perkalian skalar.
\begin{align*}
  f_2\bigl(r\cdot(c_1x+c_2y+c_3z)\bigr)
  &=f_2(rc_1x+rc_2y+rc_3z)   \\
  &=rc_2+rc_3x+rc_1x^2   \\
  &=r\cdot(c_2+c_3x+c_1x^2)   \\
  &=r\cdot\, f_2(c_1x+c_2y+c_3z)
\end{align*}
Jadi $f_2$ merupakan isomorfisme.
Kita menulis $V\isomorphicto\polyspace_2$.
\end{example}

\begin{example}
Setiap ruang isomorfik dengan dirinya sendiri di bawah pemetaan identitas.
Pemeriksaannya mudah.
\end{example}

\begin{definition} \label{df:Automorphism}
%<*df:Automorphism>
Sebuah \definend{automorfisme}\index{automorfisme}
adalah isomorfisme suatu ruang dengan
dirinya sendiri\index{isomorfisme!suatu ruang dengan dirinya sendiri}.
%</df:Automorphism>
\end{definition}

\begin{example} \label{exam:RigidPlaneMapsAutos}
%<*ex:RigidPlaneMapsAutos0>
Sebuah pemetaan \definend{dilatasi}\index{dilatasi}%
\index{automorfisme!dilatasi}
$\map{d_s}{\Re^2}{\Re^2}$ yang mengalikan semua vektor dengan skalar tak nol
$s$ merupakan automorfisme dari $\Re^2$.
%</ex:RigidPlaneMapsAutos0>
\begin{center} \small
  \includegraphics{map/mp/ch3.14}                                
\end{center}
%<*ex:RigidPlaneMapsAutos1>
Automorfisme lainnya adalah \definend{rotasi\/}\index{rotasi}%
\index{automorfisme!rotasi}
atau \definend{pemetaan perputaran}\index{pemetaan perputaran},
$\map{t_{\theta}}{\Re^2}{\Re^2}$, yang memutar semua vektor melalui
sudut~$\theta$.
%</ex:RigidPlaneMapsAutos1>
\begin{center} \small
  \includegraphics{map/mp/ch3.15}                                
\end{center}
%<*ex:RigidPlaneMapsAutos2>
Jenis ketiga automorfisme dari $\Re^2$ adalah pemetaan
$\map{f_\ell}{\Re^2}{\Re^2}$ yang \definend{membalik}
atau \definend{mencerminkan}\index{pencerminan terhadap garis}%
\index{automorfisme!pencerminan}
semua vektor terhadap suatu garis $\ell$ yang melalui titik asal.
%</ex:RigidPlaneMapsAutos2>
\begin{center} \small
  \includegraphics{map/mp/ch3.16}                                
\end{center}
Pemeriksaan bahwa semuanya merupakan automorfisme terdapat pada
\nearbyexercise{exer:RigidPlaneMapsAutos}.
\end{example}

\begin{example} \label{ex:OORThreeToFourLinCom}
Tinjau ruang $\polyspace_5$ yang terdiri atas polinomial berderajat~$5$
atau kurang, serta pemetaan $f$ yang membawa polinomial $p(x)$ ke $p(x-1)$.
Sebagai contoh, di bawah pemetaan ini $x^2\mapsto (x-1)^2=x^2-2x+1$
dan $x^3+2x\mapsto (x-1)^3+2(x-1)=x^3-3x^2+5x-3$.
Pemetaan ini merupakan automorfisme ruang tersebut;
pemeriksaannya terdapat pada \nearbyexercise{exer:PolyToPolyLinSubst}.

Isomorfisme $\polyspace_5$ dengan dirinya sendiri ini tidak sekadar memberi
tahu kita bahwa ruang tersebut ``sama'' dengan dirinya sendiri.
Isomorfisme ini juga memberi kita wawasan tentang struktur ruang tersebut.
Di bawah ini terdapat suatu keluarga parabola, yaitu grafik dari
unsur-unsur $\polyspace_5$.
Masing-masing memiliki titik puncak pada $y=-1$; parabola paling kiri
memiliki akar pada $-2.25$ dan $-1.75$, parabola berikutnya memiliki akar
pada $-1.25$ dan $-0.75$, dan seterusnya.
\begin{center}
  \includegraphics{map/mp/ch3.13}                                
\end{center}
Penggantian \( x \) dengan \( x-1 \) dalam argumen sebarang fungsi
menggeser grafiknya satu satuan ke kanan.
Jadi, $f(p_0)=p_1$, dan tindakan $f$ adalah menggeser semua parabola
satu satuan ke kanan.
Perhatikan bahwa gambar sebelum $f$ diterapkan sama dengan gambar setelah
$f$ diterapkan: ketika setiap parabola bergerak ke kanan, parabola lain
masuk dari kiri untuk menggantikannya.
Hal yang sama juga berlaku untuk grafik kubik, dan sebagainya.
Jadi, automorfisme $f$ menyatakan gagasan bahwa $P_5$ memiliki suatu
homogenitas horizontal:
jika kita menggambar dua gambar yang menampilkan semua unsur $\polyspace_5$,
satu berpusat pada $x=0$ dan yang lain berpusat pada~$x=1$, kedua gambar
tersebut tidak dapat dibedakan.
\end{example}

Seperti diuraikan pada pembukaan bagian ini, setelah memberikan definisi
isomorfisme, selanjutnya kita akan mendukung tesis bahwa definisi tersebut
menangkap intuisi kita mengenai ruang-ruang vektor yang sama.
Pertama, definisinya sendiri meyakinkan:~sebuah ruang vektor terdiri atas
suatu himpunan dan suatu struktur, dan definisi itu hanya mensyaratkan bahwa
himpunan-himpunannya berkorespondensi serta strukturnya juga berkorespondensi.
Contoh-contoh di atas juga meyakinkan, seperti
\nearbyexample{exam:TwoWideIsoTwoTall}, yang memperlihatkan bahwa ruang-ruang
isomorfik sama dalam segala hal yang relevan.
Kadang-kadang, jika \( V\isomorphicto W \), orang mengatakan bahwa
``\( W \) hanyalah \( V \) yang dicat hijau''\Dash perbedaannya semata-mata
bersifat kosmetik.

Hasil-hasil di bawah ini semakin mendukung pendapat kita bahwa di bawah
suatu isomorfisme, semua hal yang penting pada kedua ruang vektor saling
berkorespondensi.
Karena kita memperkenalkan ruang vektor untuk mempelajari kombinasi linear,
``penting'' berarti ``berkaitan dengan kombinasi linear.''
Yang tidak penting adalah cara vektor-vektor itu ditampilkan secara tipografis
(atau warnanya!\spacefactor1000).

\begin{lemma}       \label{le:IsoSendsZeroToZero}
%<*lm:IsoSendsZeroToZero>
Sebuah isomorfisme memetakan vektor nol ke vektor nol.
%</lm:IsoSendsZeroToZero>
\end{lemma}

\begin{proof}
%<*pf:IsoSendsZeroToZero>
Jika \( \map{f}{V}{W} \) adalah isomorfisme, tetapkan suatu
\( \vec{v}\in V \).
Maka \( f(\zero_V)=f(0\cdot\vec{v})=0\cdot f(\vec{v})=\zero_W \). 
%</pf:IsoSendsZeroToZero>
\end{proof}

% The definition of isomorphism requires that sums of two vectors correspond and
% that so do scalar multiples.
% We can extend that to say that all linear combinations correspond.

\begin{lemma}       \label{le:PresStructIffPresCombos}
%<*lm:PresStructIffPresCombos>
Untuk sebarang pemetaan \( \map{f}{V}{W} \) antara ruang-ruang vektor,
pernyataan-pernyataan berikut ekuivalen.
\begin{tfae}
   \item \( f \) mempertahankan struktur\index{fungsi!mempertahankan struktur}
     \begin{equation*}
        f(\vec{v}_1+\vec{v}_2)=f(\vec{v}_1)+f(\vec{v}_2)
         \quad\text{dan}\quad
        f(c\vec{v})=c\,f(\vec{v})
     \end{equation*}
   \item \( f \) mempertahankan kombinasi linear dua vektor
     \begin{equation*}
        f(c_1\vec{v}_1+c_2\vec{v}_2)=c_1f(\vec{v}_1)+c_2f(\vec{v}_2)
     \end{equation*}
   \item \( f \) mempertahankan kombinasi linear sebarang banyak vektor
     yang berhingga
     \begin{equation*}
        f(c_1\vec{v}_1+\dots+c_n\vec{v}_n)=
           c_1f(\vec{v}_1)+\dots+c_nf(\vec{v}_n)
     \end{equation*}
\end{tfae}
%</lm:PresStructIffPresCombos>
\end{lemma}

\begin{proof}
%<*pf:PresStructIffPresCombos0>
Karena implikasi \mbox{$\text{(3)}\!\implies\!\text{(2)}$} dan 
\mbox{$\text{(2)}\!\implies\!\text{(1)}$}
sudah jelas, kita hanya perlu menunjukkan bahwa
\mbox{$\text{(1)}\!\implies\!\text{(3)}$}.
Maka, andaikan pernyataan~(1).
Kita akan membuktikan~(3) dengan induksi pada banyaknya suku $n$.
%</pf:PresStructIffPresCombos0>

%<*pf:PresStructIffPresCombos1>
Kasus dasar satu suku, yaitu
\( f(c\vec{v}_1)=c\,f(\vec{v}_1) \), sudah tercakup dalam klausa kedua
pernyataan~(1).
%</pf:PresStructIffPresCombos1>

%<*pf:PresStructIffPresCombos2>
Untuk langkah induksi, andaikan pernyataan~(3) berlaku setiap kali terdapat
\( k \) suku atau kurang.
% , that is, whenever
% $n=1$, or $n=2$, \ldots, or $n=k$.
Tinjau kasus dengan $k+1$ suku.
Gunakan bagian pertama dari~(1) untuk memisahkan jumlah pada tanda `$+$'
terakhir.
\begin{equation*}
  f(c_1\vec{v}_1+\dots+c_k\vec{v}_k+c_{k+1}\vec{v}_{k+1})
  =f(c_1\vec{v}_1+\dots+c_k\vec{v}_k)+f(c_{k+1}\vec{v}_{k+1})
\end{equation*}
Gunakan hipotesis induksi untuk memisahkan jumlah $k$ suku di sebelah kiri.
\begin{equation*}
  =f(c_1\vec{v}_1)+\dots+f(c_k\vec{v}_k)+f(c_{k+1}\vec{v}_{k+1})
\end{equation*}
Sekarang bagian kedua dari~(1) memberikan
\begin{equation*}
  =c_1\,f(\vec{v}_1)+\dots+c_k\,f(\vec{v}_k)+c_{k+1}\,f(\vec{v}_{k+1})
\end{equation*}
ketika diterapkan sebanyak $k+1$~kali.
%</pf:PresStructIffPresCombos2>
\end{proof}

% In addition to adding to the intuition that the definition of isomorphism
% does indeed preserve the things of interest in a vector space, 
\noindent Kita sering menggunakan butir~(2) untuk menyederhanakan verifikasi
bahwa suatu pemetaan mempertahankan struktur.

Sebagai penutup, berikut ringkasannya.
Dalam bab sebelumnya, setelah memberikan definisi ruang vektor,
kita melihat berbagai contoh dan mencatat bahwa beberapa ruang tampaknya
pada hakikatnya sama dengan ruang lain.
Di sini kita telah mendefinisikan relasi `\( \cong \)' dan berargumen bahwa
relasi tersebut merupakan cara yang tepat untuk menyatakan secara cermat
apa yang kita maksud dengan ``sama'', karena relasi itu mempertahankan
ciri-ciri penting suatu ruang vektor\Dash khususnya, kombinasi linear.
Pada bagian berikutnya kita akan menunjukkan bahwa isomorfisme merupakan
relasi ekuivalensi sehingga mempartisi kumpulan ruang vektor.


\begin{exercises}
  \recommended \item \label{exer:PThreeIsoRFour}
   Dengan menggunakan \nearbyexample {ex:CheckMapIsIso} sebagai model,
   verifikasikan bahwa kedua korespondensi yang diberikan sebelum definisi
   merupakan isomorfisme.
   \begin{exparts*}
      \partsitem \nearbyexample{exam:TwoWideIsoTwoTall}
      \partsitem \nearbyexample{exam:PolyTwoIsoRThree}
    \end{exparts*}
    \begin{answer}
     \begin{exparts}
     \partsitem Namakan pemetaan tersebut $f$.
       \begin{equation*}
         \rowvec{a &b} \mapsunder{f} \colvec{a \\ b}
       \end{equation*}
       Pemetaan ini injektif karena jika $f$ membawa dua unsur domain ke
       bayangan yang sama, yakni jika
       $f\left(\rowvec{a &b}\right)=f\left(\rowvec{c &d}\right)$,
       maka definisi $f$ memberikan
       \begin{equation*}
          \colvec{a \\ b}=\colvec{c \\ d}
       \end{equation*}
       dan karena vektor-vektor kolom sama hanya jika komponen-komponennya sama,
       kita memperoleh $a=c$ dan $b=d$.
       Jadi, jika $f$ memetakan dua vektor baris dari domain ke vektor kolom
       yang sama, maka kedua vektor baris itu sama:
       $\rowvec{a &b}=\rowvec{c &d}$.

       Untuk menunjukkan bahwa $f$ surjektif, kita harus menunjukkan bahwa
       setiap unsur kodomain $\Re^2$ adalah bayangan suatu vektor baris di
       bawah $f$. Hal ini mudah;
       \begin{equation*}
          \colvec{x \\ y} 
       \end{equation*}
       adalah $f\left(\rowvec{x &y}\right)$.

       Perhitungan yang menunjukkan dipertahankannya penjumlahan adalah
       sebagai berikut.
       \begin{multline*}
         f\left(\rowvec{a &b}+\rowvec{c &d}\right) 
           = f\left(\rowvec{a+c &b+d}\right)           
           = \colvec{a+c \\ b+d}                               \\
           = \colvec{a \\ b}+\colvec{c \\ d}         
           = f\left(\rowvec{a &b}\right)+f\left(\rowvec{c &d}\right)
       \end{multline*}
       Perhitungan untuk perkalian skalar serupa.
       \begin{equation*}
         f\left(r\cdot\rowvec{a &b}\right) 
           = f\left(\rowvec{ra &rb}\right)          
           = \colvec{ra \\ rb}          
           = r\cdot\colvec{a \\ b}          
           = r\cdot f\left(\rowvec{a &b}\right)
       \end{equation*}
     \partsitem Nyatakan pemetaan dari \nearbyexample{exam:PolyTwoIsoRThree}
      dengan $f$.
      Untuk menunjukkan bahwa pemetaan tersebut injektif,
      andaikan \( f(a_0+a_1x+a_2x^2)=f(b_0+b_1x+b_2x^2) \).
      Maka, berdasarkan definisi fungsi,
      \begin{equation*}
        \colvec{a_0 \\ a_1 \\ a_2}
        =\colvec{b_0 \\ b_1 \\ b_2}
      \end{equation*}
      sehingga \( a_0=b_0 \), \( a_1=b_1 \), dan \( a_2=b_2 \).
      Jadi \( a_0+a_1x+a_2x^2=b_0+b_1x+b_2x^2 \), sehingga \( f \)
      injektif.

      Fungsi $f$ surjektif karena terdapat polinomial yang dipetakan ke
      \begin{equation*}
        \colvec{a \\ b \\ c}
      \end{equation*}
      oleh \( f \), yaitu \( a+bx+cx^2 \).

      Mengenai struktur, perhitungan berikut menunjukkan bahwa $f$
      mempertahankan penjumlahan
      \begin{multline*}
        f\left(\,(a_0+a_1x+a_2x^2)+(b_0+b_1x+b_2x^2)\,\right)           \\
        \begin{aligned}
        &=f\left(\,(a_0+b_0)+(a_1+b_1)x+(a_2+b_2)x^2\,\right)       \\
        &=\colvec{a_0+b_0 \\ a_1+b_1 \\ a_2+b_2}                    \\
        &=\colvec{a_0 \\ a_1 \\ a_2} +\colvec{b_0 \\ b_1 \\ b_2}   \\
        &=f(a_0+a_1x+a_2x^2)+f(b_0+b_1x+b_2x^2)
        \end{aligned}
      \end{multline*}
      dan perhitungan berikut menunjukkan
      \begin{align*}
        f(\,r(a_0+a_1x+a_2x^2)\,)
        &=f(\,(ra_0)+(ra_1)x+(ra_2)x^2\,)    \\
        &=\colvec{ra_0 \\ ra_1 \\ ra_2}     \\
        &=r\cdot\colvec{a_0 \\ a_1 \\ a_2}     \\
        &=r\,f(a_0+a_1x+a_2x^2)
      \end{align*}   
      bahwa pemetaan tersebut mempertahankan perkalian skalar.
    \end{exparts} 
   \end{answer}
  \recommended \item 
    Untuk pemetaan
    \( \map{f}{\polyspace_1}{\Re^2} \) yang diberikan oleh
    \begin{equation*}
       a+bx\mapsunder{f}\colvec{a-b \\ b}
    \end{equation*}
    tentukan bayangan setiap unsur domain berikut.
    \begin{exparts*}
      \partsitem \( 3-2x \)
      \partsitem \( 2+2x \)
      \partsitem \( x \)
    \end{exparts*}
    Tunjukkan bahwa pemetaan ini merupakan isomorfisme.
    \begin{answer}
       Berikut adalah bayangan-bayangannya.
       \begin{exparts*}
        \partsitem \( \colvec[r]{5 \\ -2} \)
        \partsitem \( \colvec[r]{0 \\ 2} \)
        \partsitem \( \colvec[r]{-1 \\ 1} \)
      \end{exparts*}

      Untuk membuktikan bahwa $f$ injektif, andaikan bahwa pemetaan ini membawa
      dua polinomial linear ke bayangan yang sama,
      $f(a_1+b_1x)=f(a_2+b_2x)$. Maka
      \begin{equation*}
        \colvec{a_1-b_1 \\ b_1}=\colvec{a_2-b_2 \\ b_2}
      \end{equation*}
      dan karena vektor-vektor kolom sama hanya jika komponen-komponennya sama,
      diperoleh $b_1=b_2$ dan $a_1=a_2$.
      Hal ini menunjukkan bahwa kedua polinomial linear itu sama, sehingga
      $f$ injektif.

      Untuk menunjukkan bahwa $f$ surjektif, perhatikan bahwa unsur kodomain
      \begin{equation*}
         \colvec{s \\ t}
      \end{equation*}
      ini merupakan bayangan unsur domain $(s+t)+tx$.

      Untuk memeriksa bahwa $f$ mempertahankan struktur,
      kita dapat menggunakan butir~(2) dari
      \nearbylemma{le:PresStructIffPresCombos}.
      \begin{align*}
        f\left(c_1\cdot (a_1+b_1x)+c_2\cdot (a_2+b_2x)\right)
        &=f\left((c_1a_1+c_2a_2)+(c_1b_1+c_2b_2)x\right)                \\
        &=\colvec{(c_1a_1+c_2a_2)-(c_1b_1+c_2b_2) \\ c_1b_1+c_2b_2}       \\  
        &=c_1\cdot\colvec{a_1-b_1 \\ b_1}+c_2\cdot\colvec{a_2-b_2 \\ b_2}  \\
        &=c_1\cdot f(a_1+b_1x)+c_2\cdot f(a_2+b_2x)
      \end{align*}
     \end{answer}
  \item \label{exer:NatMapAlsoIso}
    Tunjukkan bahwa pemetaan alami $f_1$ dari
    \nearbyexample{ex:LinComboThreeIsoPTwo} 
    merupakan isomorfisme.
    \begin{answer}
      Untuk memverifikasi bahwa pemetaan tersebut injektif, andaikan
       $f_1(c_1x+c_2y+c_3z)=f_1(d_1x+d_2y+d_3z)$.
      Maka $c_1+c_2x+c_3x^2=d_1+d_2x+d_3x^2$ berdasarkan definisi $f_1$.
      Unsur-unsur $\polyspace_2$ sama hanya jika koefisien-koefisiennya sama,
      sehingga hal ini menyiratkan $c_1=d_1$, $c_2=d_2$, dan $c_3=d_3$.
      Oleh karena itu, $f_1(c_1x+c_2y+c_3z)=f_1(d_1x+d_2y+d_3z)$ menyiratkan
      $c_1x+c_2y+c_3z=d_1x+d_2y+d_3z$, sehingga $f_1$ injektif.

      Untuk memverifikasi bahwa pemetaan tersebut surjektif, tinjau sebarang
      unsur kodomain $a_1+a_2x+a_3x^2$ dan perhatikan bahwa unsur itu memang
      merupakan bayangan suatu unsur domain, yaitu
      $f_1(a_1x+a_2y+a_3z)$.
      (Sebagai contoh, $0+3x+6x^2=f_1(0x+3y+6z)$.)

      Perhitungan yang memeriksa bahwa $f_1$ mempertahankan penjumlahan adalah
      sebagai berikut.
      \begin{multline*}
        f_1\left(\,(c_1x+c_2y+c_3z)+(d_1x+d_2y+d_3z)\,\right)            \\
          \begin{aligned}
          &=f_1\left(\,(c_1+d_1)x+(c_2+d_2)y+(c_3+d_3)z\,\right)  \\
          &=(c_1+d_1)+(c_2+d_2)x+(c_3+d_3)x^2  \\
          &=(c_1+c_2x+c_3x^2)+(d_1+d_2x+d_3x^2)  \\
          &=f_1(c_1x+c_2y+c_3z)+f_1(d_1x+d_2y+d_3z)
          \end{aligned}
      \end{multline*}

      Pemeriksaan bahwa $f_1$ mempertahankan perkalian skalar adalah sebagai
      berikut.
      \begin{align*}
        f_1(\,r\cdot (c_1x+c_2y+c_3z)\,)
          &=f_1(\,(rc_1)x+(rc_2)y+(rc_3)z\,)  \\
          &=(rc_1)+(rc_2)x+(rc_3)x^2  \\
          &=r\cdot (c_1+c_2x+c_3x^2)  \\
          &=r\cdot f_1(c_1x+c_2y+c_3z)
      \end{align*}
    \end{answer}
  \item Tunjukkan bahwa pemetaan $\map{t}{\polyspace_2}{\polyspace_2}$ yang
    diberikan oleh $t(ax^2+bx+c)=bx^2-(a+c)x+a$ merupakan isomorfisme.
    \begin{answer}
      Untuk melihat bahwa pemetaan tersebut injektif, andaikan
      $t(\vec{v}_1)=t(\vec{v}_2)$ dengan tujuan menyimpulkan
      $\vec{v}_1=\vec{v}_2$.
      Artinya, $t(a_1x^2+b_1x+c_1)=t(a_2x^2+b_2x+c_2)$.
      Maka $b_1x^2-(a_1+c_1)x+a_1=b_2x^2-(a_2+c_2)x+a_2$.
      Karena polinomial kuadrat sama hanya jika suku kuadrat, suku konstan,
      dan suku linearnya sama, kita menyimpulkan bahwa $b_1=b_2$,
      $a_1=a_2$, dan dari sini $c_1=c_2$.
      Oleh karena itu, $a_1x^2+b_1x+c_1=a_2x^2+b_2x+c_2$ dan fungsi tersebut
      injektif.

      Untuk melihat bahwa pemetaan tersebut surjektif, andaikan diberikan
      suatu unsur~$\vec{w}$ dari kodomain, lalu carilah unsur~$\vec{v}$ dari
      domain yang dipetakan kepadanya.
      Misalkan unsur kodomain itu~$\vec{w}=px^2+qx+r$.
      Perhatikan bahwa jika $\vec{v}=rx^2+px+(-q-r)$, maka
      $t(\vec{v})=\vec{w}$. Jadi~$t$ surjektif.

      Untuk melihat bahwa pemetaan tersebut merupakan homomorfisme, kita
      menunjukkan bahwa pemetaan itu menghormati kombinasi linear dua unsur.
      Berdasarkan \nearbylemma{le:PresStructIffPresCombos}, hal ini akan
      menunjukkan bahwa pemetaan tersebut mempertahankan operasi-operasinya.
      \begin{multline*}
        t(r_1(a_1x^2+b_1x+c_1)+r_2(a_2x^2+b_2x+c_2))              \\ 
        \begin{split} \quad 
        &=t((r_1a_1+r_2a_2)x^2+(r_1b_1+r_2b_2)x+(r_1c_1+r_2c_2))   \\
        &=(r_1b_1+r_2b_2)x^2-((r_1a_1+r_2a_2)+(r_1c_1+r_2c_2))x+(r_1a_1+r_2a_2)  \\
        &=(r_1b_1)x^2-(r_1a_1+r_1c_1)x+r_1a_1
           +(r_2b_2)x^2-(r_2a_2+r_2c_2)x+r_2a_2                       \\
        &=r_1t(a_1x^2+b_1x+c_1)+r_2t(a_2x^2+b_2x+c_2)
        \end{split}
      \end{multline*}      
    \end{answer}
  \recommended \item Verifikasikan bahwa pemetaan berikut merupakan isomorfisme:
    $\map{h}{\Re^4}{\matspace_{\nbyn{2}}}$ yang diberikan oleh
    \begin{equation*}
      \colvec{a \\ b \\ c \\ d}
      \mapsto
      \begin{mat}
        c  &a+d \\
        b  &d
      \end{mat}
    \end{equation*}
    \begin{answer}
      Pertama-tama kita memverifikasi bahwa $h$ injektif.
      Untuk itu, kita akan menunjukkan bahwa $h(\vec{v}_1)=h(\vec{v}_2)$
      menyiratkan $\vec{v}_1=\vec{v}_2$.
      Maka andaikan
      \begin{equation*}
          h(\vec{v}_1)
          =
          h(\colvec{a_1 \\ b_1 \\ c_1 \\ d_1})
          =
          h(\colvec{a_2 \\ b_2 \\ c_2 \\ d_2})
          =h(\vec{v}_2)
      \end{equation*}
      yang memberikan
      \begin{equation*}
          \begin{mat}
            c_1  &a_1+d_1 \\
            b_1  &d_1
          \end{mat}
          =
          \begin{mat}
            c_2  &a_2+d_2 \\
            b_2  &d_2
          \end{mat}
      \end{equation*}
      dari sini kita menyimpulkan bahwa
      $c_1=c_2$ (dari entri kiri-atas),
      $b_1=b_2$ (dari entri kiri-bawah),
      $d_1=d_2$ (dari entri kanan-bawah),
      dan dari kesamaan terakhir ini kita memperoleh $a_1=a_2$
      (dari entri kanan-atas).
      Oleh karena itu, $\vec{v}_1=\vec{v}_2$.

      Selanjutnya kita akan menunjukkan bahwa pemetaan tersebut surjektif,
      yakni bahwa setiap unsur kodomain $\matspace_{\nbyn{2}}$ merupakan
      bayangan suatu vektor kolom berukuran empat dari domain.
      Maka, jika diberikan
      \begin{equation*}
          \vec{w}=
          \begin{mat}
            m  &n \\
            p  &q
          \end{mat}
          \in\matspace_{\nbyn{2}}
      \end{equation*}
      perhatikan bahwa matriks tersebut merupakan bayangan vektor domain berikut.
      \begin{equation*}
        \vec{v}=
        \colvec{n-q \\ p \\ m \\ q}
      \end{equation*}

      Sebagai penutup, kita memverifikasi bahwa pemetaan tersebut
      mempertahankan kombinasi linear.
      Berdasarkan \nearbylemma{le:PresStructIffPresCombos}, hal ini menunjukkan
      bahwa pemetaan tersebut mempertahankan operasi-operasinya.
      \begin{align*}
          h(r_1\cdot\colvec{a_1 \\ b_1 \\ c_1 \\ d_1}
            +
            r_2\cdot\colvec{a_2 \\ b_2 \\ c_2 \\ d_2})
          &=h(\colvec{r_1a_1+r_2a_2 \\ r_1b_1+r_2b_2 \\ r_1c_1+r_2c_2 \\ r_1d_1+r_2d_2})  \\
          &=
          \begin{mat}
             r_1c_1+r_2c_2 &(r_1a_1+r_2a_2)+(r_1d_1+r_2d_2)  \\
             r_1b_1+r_2b_2 &r_1d_1+r_2d_2
          \end{mat}                               \\
          &=
          r_1\begin{mat}
             c_1 &a_1+d_1  \\
             b_1 &d_1
          \end{mat}                               
          +
          r_2\begin{mat}
             c_2 &a_2+d_2  \\
             b_2 &d_2
          \end{mat}                               \\
         &=r_1\cdot h(\colvec{a_1 \\ b_1 \\ c_1 \\ d_1})
            +
            r_2\cdot h(\colvec{a_2 \\ b_2 \\ c_2 \\ d_2})
      \end{align*}      
    \end{answer}
  \recommended \item 
   Tentukan apakah setiap pemetaan berikut merupakan isomorfisme.
   Jika ya, buktikan; jika tidak, nyatakan syarat yang gagal dipenuhinya.
    \begin{exparts}
      \partsitem \( \map{f}{\matspace_{\nbyn{2}}}{\Re} \) yang diberikan oleh
        \begin{equation*}
           \begin{mat}
             a  &b  \\
             c  &d
           \end{mat}
           \mapsto
           ad-bc
        \end{equation*}
      \partsitem \( \map{f}{\matspace_{\nbyn{2}}}{\Re^4} \) yang diberikan oleh
        \begin{equation*}
           \begin{mat}
             a  &b  \\
             c  &d
           \end{mat}
           \mapsto
           \colvec{a+b+c+d \\ a+b+c \\ a+b \\ a}
        \end{equation*}
      \partsitem \( \map{f}{\matspace_{\nbyn{2}}}{\polyspace_3} \) yang diberikan oleh
        \begin{equation*}
           \begin{mat}
             a  &b  \\
             c  &d
           \end{mat}
           \mapsto
           c+(d+c)x+(b+a)x^2+ax^3
        \end{equation*}
      \partsitem \( \map{f}{\matspace_{\nbyn{2}}}{\polyspace_3} \) yang diberikan oleh
        \begin{equation*}
           \begin{mat}
             a  &b  \\
             c  &d
           \end{mat}
           \mapsto
           c+(d+c)x+(b+a+1)x^2+ax^3
        \end{equation*}
    \end{exparts}
  \begin{answer}
    \begin{exparts}
      \partsitem Tidak; pemetaan ini tidak injektif.
        Secara khusus, matriks yang semua entrinya nol dipetakan ke bayangan
        yang sama dengan matriks yang semua entrinya satu.
      \partsitem Ya, pemetaan ini merupakan isomorfisme.

        Pemetaan ini injektif:
        \begin{align*}
          &\text{jika }
          f(\begin{mat}
              a_1  &b_1  \\
              c_1  &d_1
            \end{mat})
         =f(\begin{mat}
              a_2  &b_2  \\
              c_2  &d_2
            \end{mat})                                     \\
          &\qquad\text{ maka }
          \colvec{a_1+b_1+c_1+d_1 \\ a_1+b_1+c_1 \\ a_1+b_1 \\ a_1}
          =\colvec{a_2+b_2+c_2+d_2 \\ a_2+b_2+c_2 \\ a_2+b_2 \\ a_2}
        \end{align*}
        memberikan \( a_1=a_2 \), \( b_1=b_2 \),
        \( c_1=c_2 \), dan \( d_1=d_2 \).

        Pemetaan ini surjektif, karena persamaan berikut menunjukkan
        \begin{equation*}
          \colvec{x \\ y \\ z \\ w}
         =f(\begin{mat}
              w    &z-w  \\
              y-z  &x-y
            \end{mat})
        \end{equation*}
        bahwa sebarang vektor kolom berukuran empat merupakan bayangan suatu
        matriks $\nbyn{2}$.

        Terakhir, pemetaan tersebut mempertahankan kombinasi linear
        \begin{multline*}
          f(\,r_1\cdot\begin{mat}
              a_1  &b_1  \\
              c_1  &d_1
            \end{mat}
            +r_2\cdot\begin{mat}
              a_2  &b_2  \\
              c_2  &d_2
            \end{mat}\,)                                   \\
          \begin{aligned}
          &=f(\begin{mat}
              r_1a_1+r_2a_2  &r_1b_1+r_2b_2  \\
              r_1c_1+r_2c_2  &r_1d_1+r_2d_2
            \end{mat})                            \\
          &=\colvec{r_1a_1+\dots+r_2d_2 \\ r_1a_1+\dots+r_2c_2 \\ 
                       r_1a_1+\dots+r_2b_2 \\ r_1a_1+r_2a_2}          \\
          &=r_1\cdot\colvec{a_1+\dots+d_1 \\ a_1+\dots+c_1 \\ 
                       a_1+b_1 \\ a_1}
          +r_2\cdot\colvec{a_2+\dots+d_2 \\ a_2+\dots+c_2 \\ 
                       a_2+b_2 \\ a_2}                         \\
          &=r_1\cdot f(\begin{mat}
                  a_1  &b_1  \\
                  c_1  &d_1
            \end{mat})
          +r_2\cdot f(\begin{mat}
                  a_2  &b_2  \\
                  c_2  &d_2
            \end{mat})
         \end{aligned}
        \end{multline*}
        sehingga butir~(2) dari \nearbylemma{le:PresStructIffPresCombos}
        menunjukkan bahwa pemetaan itu mempertahankan struktur.
      \partsitem Ya, pemetaan ini merupakan isomorfisme.

        Untuk menunjukkan bahwa pemetaan ini injektif, andaikan dua unsur
        domain memiliki bayangan yang sama di bawah $f$.
        \begin{equation*}
          f(\begin{mat}
              a_1  &b_1  \\
              c_1  &d_1
            \end{mat})
          =f(\begin{mat}
              a_2  &b_2  \\
              c_2  &d_2
            \end{mat})
        \end{equation*}
        Berdasarkan definisi $f$, hal ini memberikan
        $c_1+(d_1+c_1)x+(b_1+a_1)x^2+a_1x^3
          =c_2+(d_2+c_2)x+(b_2+a_2)x^2+a_2x^3$
        lalu fakta bahwa polinomial sama hanya jika koefisien-koefisiennya sama
        memberikan sistem persamaan linear
        \begin{align*}
          c_1      &=  c_2      \\
          d_1+c_1  &=  d_2+c_2  \\
          b_1+a_1  &=  b_2+a_2  \\
          a_1      &=  a_2
        \end{align*}
        yang hanya memiliki solusi $a_1=a_2$, $b_1=b_2$, $c_1=c_2$, dan
        $d_1=d_2$.

        Untuk menunjukkan bahwa $f$ surjektif, perhatikan bahwa
        $p+qx+rx^2+sx^3$ merupakan bayangan matriks berikut di bawah $f$.
        \begin{equation*}
          \begin{mat}
            s  &r-s   \\
            p  &q-p
          \end{mat}
        \end{equation*}

        Kita dapat memeriksa bahwa $f$ mempertahankan struktur dengan
        menggunakan butir~(2) dari \nearbylemma{le:PresStructIffPresCombos}.
        \begin{multline*}
          f(r_1\cdot\begin{mat}
              a_1  &b_1  \\
              c_1  &d_1
            \end{mat}
            +r_2\cdot\begin{mat}
              a_2  &b_2  \\
              c_2  &d_2
            \end{mat})                                      \\
           \begin{aligned}
           &=
           f(\begin{mat}
               r_1a_1+r_2a_2  &r_1b_1+r_2b_2  \\
               r_1c_1+r_2c_2  &r_1d_1+r_2d_2
             \end{mat})                         \\
           &=\begin{aligned}[t]
               &(r_1c_1+r_2c_2)
               +(r_1d_1+r_2d_2+r_1c_1+r_2c_2)x                 \\
               &\mbox{}\quad +(r_1b_1+r_2b_2+r_1a_1+r_2a_2)x^2
                 +(r_1a_1+r_2a_2)x^3
             \end{aligned}                   \\
           &=\begin{aligned}[t]
                &r_1\cdot\left(c_1+(d_1+c_1)x+(b_1+a_1)x^2+a_1x^3\right) \\
                &\mbox{}\quad +r_2\cdot\left(c_2+(d_2+c_2)x
                            +(b_2+a_2)x^2+a_2x^3\right)
              \end{aligned} \\
           &=r_1\cdot f(\begin{mat}
              a_1  &b_1  \\
              c_1  &d_1
            \end{mat})
            +r_2\cdot f(\begin{mat}
              a_2  &b_2  \\
              c_2  &d_2
            \end{mat})
            \end{aligned}
        \end{multline*}
      \partsitem Tidak, pemetaan ini tidak mempertahankan struktur.
        Sebagai contoh, pemetaan ini tidak membawa matriks yang semua entrinya
        nol ke polinomial nol.
    \end{exparts}  
   \end{answer}
  \item 
    Tunjukkan bahwa pemetaan \( \map{f}{\Re^1}{\Re^1} \) yang diberikan oleh
    \( f(x)=x^3 \) bersifat injektif dan surjektif.
    Apakah pemetaan tersebut merupakan isomorfisme?
    \begin{answer}
      Pemetaan ini injektif dan surjektif, sehingga merupakan korespondensi,
      karena memiliki invers (yaitu \( f^{-1}(x)=\sqrt[3]{x} \)).
      Namun, pemetaan ini bukan isomorfisme.
      Sebagai contoh, \( f(1)+f(1)\neq f(1+1) \).
    \end{answer}
  \recommended \item 
    Lihat \nearbyexample{exam:TwoWideIsoTwoTall}.
    Buat dua isomorfisme lainnya
    (tentu saja, Anda juga harus memverifikasi bahwa keduanya memenuhi
    syarat-syarat dalam definisi isomorfisme).
    \begin{answer}
     Banyak pemetaan yang mungkin.
     Berikut dua di antaranya.
     \begin{equation*}
       \rowvec{a &b}\mapsto\colvec{b \\ a}
       \quad\text{dan}\quad
       \rowvec{a &b}\mapsto\colvec{2a \\ b}
     \end{equation*} 
     Verifikasinya merupakan adaptasi langsung dari verifikasi-verifikasi di atas.
    \end{answer}
  \item 
    Lihat \nearbyexample{exam:PolyTwoIsoRThree}.
    Buat dua isomorfisme lainnya
    (dan verifikasikan bahwa keduanya memenuhi syarat-syarat tersebut).
    \begin{answer}
      Berikut dua di antaranya.
      \begin{equation*}
        a_0+a_1x+a_2x^2 \mapsto \colvec{a_1 \\ a_0 \\ a_2}
        \quad\text{dan}\quad
        a_0+a_1x+a_2x^2 \mapsto \colvec{a_0+a_1 \\ a_1 \\ a_2}
      \end{equation*}
      Verifikasinya langsung (untuk pemetaan kedua, guna menunjukkan bahwa
      pemetaan itu surjektif, perhatikan bahwa
      \begin{equation*}
        \colvec{s \\ t \\ u}
      \end{equation*}
      merupakan bayangan dari $(s-t)+tx+ux^2$).
     \end{answer}
  \recommended \item 
     Tunjukkan bahwa meskipun \( \Re^2 \) sendiri bukan subruang dari
     \( \Re^3 \), ruang ini isomorfik dengan subruang bidang \( xy \)
     dari \( \Re^3 \).
     \begin{answer}
       Ruang $\Re^2$ bukan subruang dari $\Re^3$ karena bukan subhimpunan
       dari $\Re^3$.
       Vektor-vektor kolom berukuran dua dalam $\Re^2$ bukan unsur $\Re^3$.

       Misalkan
       \( P=\set{\colvec{x \\ y \\ 0}\suchthat x,y\in\Re} \), yaitu
       subruang bidang \(xy\) dari \( \Re^3 \).
       Isomorfisme alami \( \map{\iota}{\Re^2}{P} \)
       (disebut pemetaan \definend{injeksi}) adalah sebagai berikut.
       \begin{equation*}
         \colvec{x \\ y}
           \mapsunder{\iota}
         \colvec{x \\ y \\ 0}
       \end{equation*}

       Pemetaan ini injektif karena
       \begin{equation*}
         \iota(\colvec{x_1 \\ y_1})=\iota(\colvec{x_2 \\ y_2})
         \quad\text{menyiratkan}\quad
         \colvec{x_1 \\ y_1 \\ 0}=\colvec{x_2 \\ y_2 \\ 0}
       \end{equation*}
       yang selanjutnya menyiratkan $x_1=x_2$ dan $y_1=y_2$, sehingga kedua
       vektor kolom berukuran dua semula sama.

       Karena
       \begin{equation*}
         \colvec{x \\ y \\ 0}
         =\iota(\colvec{x \\ y})
       \end{equation*}
       pemetaan ini surjektif ke bidang $xy$.

       Untuk menunjukkan bahwa pemetaan ini mempertahankan struktur, kita akan
       menggunakan butir~(2) dari \nearbylemma{le:PresStructIffPresCombos}
       dan menunjukkan
       \begin{multline*}
         \iota(c_1\cdot \colvec{x_1 \\ y_1}+c_2\cdot \colvec{x_2 \\ y_2})
          = \iota(\colvec{c_1x_1+c_2x_2 \\ c_1y_1+c_2y_2})  
          =\colvec{c_1x_1+c_2x_2 \\ c_1y_1+c_2y_2 \\ 0}                     \\
          =c_1\cdot \colvec{x_1 \\ y_1 \\ 0}+c_2\cdot \colvec{x_2 \\ y_2 \\ 0}
          =c_1\cdot \iota(\colvec{x_1 \\ y_1})
             +c_2\cdot \iota(\colvec{x_2 \\ y_2})
       \end{multline*}
       bahwa pemetaan ini mempertahankan kombinasi linear dua vektor.
     \end{answer}
  \item 
    Temukan dua isomorfisme antara \( \Re^{16} \) dan
    \( \matspace_{\nbyn{4}} \).
     \begin{answer}
        Berikut dua di antaranya:
        \begin{equation*}
           \colvec{r_1 \\ r_2 \\ \vdots \\ r_{16}}
              \mapsto
           \begin{mat}
              r_1  &r_2  &\ldots  \\
                   &              \\
                   &     &\ldots  &r_{16}
           \end{mat}
           \quad\text{dan}\quad
           \colvec{r_1 \\ r_2 \\ \vdots \\ r_{16}}
              \mapsto
           \begin{mat}
              r_1    &                    \\
              r_2    &                    \\
              \vdots &   &        &\vdots \\
                   &     &        &r_{16}
           \end{mat}
        \end{equation*}     
        Verifikasi bahwa masing-masing merupakan isomorfisme mudah dilakukan.
    \end{answer}
  \recommended \item
    Untuk nilai \( k \) berapakah \( \matspace_{\nbym{m}{n}} \) isomorfik
    dengan \( \Re^{k} \)?
     \begin{answer}
        Jika $k$ adalah hasil kali \( k=mn \), berikut suatu isomorfisme.
        \begin{equation*}
           \begin{mat}
              r_1  &r_2  &\ldots  \\
                   &\vdots        \\
                   &     &\ldots  &r_{m\cdot n}
           \end{mat}
           \mapsto
           \colvec{r_1 \\ r_2 \\ \vdots \\ r_{m\cdot n}}
        \end{equation*}
        Pemeriksaan bahwa pemetaan ini merupakan isomorfisme mudah dilakukan.
      \end{answer}
  \item 
     Untuk nilai \( k \) berapakah \( \polyspace_k \) isomorfik dengan
     \( \Re^n \)?
     \begin{answer}
        Jika \( n\geq 1 \), maka \( \polyspace_{n-1}\isomorphicto\Re^n \).
        (Jika kita mengambil \( \polyspace_{-1} \) dan \( \Re^0 \) sebagai
        ruang vektor trivial, hubungan tersebut meluas satu dimensi lebih
        rendah.)
        Isomorfisme alami di antara keduanya adalah sebagai berikut.
        \begin{equation*}
          a_0+a_1x+\dots+a_{n-1}x^{n-1}
          \mapsto\colvec{a_0 \\ a_1 \\ \vdots \\ a_{n-1}}
        \end{equation*}
        Pemeriksaan bahwa pemetaan tersebut merupakan isomorfisme berlangsung
        langsung.
     \end{answer}
  \item \label{exer:PolyToPolyLinSubst}
    Buktikan bahwa pemetaan dalam \nearbyexample{ex:OORThreeToFourLinCom},
    dari \( \polyspace_5 \) ke \( \polyspace_5 \) yang diberikan oleh
    \( p(x)\mapsto p(x-1) \), merupakan isomorfisme ruang vektor.
    \begin{answer}
      Berikut pemetaan tersebut setelah dijabarkan.
      \begin{multline*}
        f(a_0+a_1x+a_2x^2+a_3x^3+a_4x^4+a_5x^5)                        \\
        \begin{aligned}
        &=a_0+a_1(x-1)+a_2(x-1)^2+a_3(x-1)^3               \\
        &\mbox{}\quad +a_4(x-1)^4+a_5(x-1)^5               \\
        &=a_0+a_1(x-1)+a_2(x^2-2x+1)         \\
        &\mbox{}\quad +a_3(x^3-3x^2+3x-1)   \\
        &\mbox{}\quad +a_4(x^4-4x^3+6x^2-4x+1) \\
        &\mbox{}\quad +a_5(x^5-5x^4+10x^3-10x^2+5x-1)       \\
        &=(a_0-a_1+a_2-a_3+a_4-a_5)  \\
        &\mbox{}\quad +(a_1-2a_2+3a_3-4a_4+5a_5)x  \\
        &\mbox{}\quad +(a_2-3a_3+6a_4-10a_5)x^2  \\
        &\quad+(a_3-4a_4+10a_5)x^3         \\
        &\mbox{}\quad +(a_4-5a_5)x^4
                        +a_5x^5
        \end{aligned}
      \end{multline*}
      Pemetaan ini merupakan korespondensi karena memiliki invers, yaitu
      pemetaan \( p(x)\mapsto p(x+1) \).

      Untuk menyelesaikan pemeriksaan bahwa pemetaan ini merupakan isomorfisme,
      kita menerapkan butir~(2) dari
      \nearbylemma{le:PresStructIffPresCombos} dan menunjukkan bahwa pemetaan
      tersebut mempertahankan kombinasi linear dua polinomial.
      Secara ringkas,
      $f(c\cdot (a_0+a_1x+\dots +a_5x^5)+d\cdot (b_0+b_1x+\dots +b_5x^5))$
      sama dengan
      \begin{multline*}
         (ca_0-ca_1+ca_2-ca_3+ca_4-ca_5+db_0-db_1+db_2-db_3+db_4-db_5)  \\
          +\dots+(ca_5+db_5)x^5                                
      \end{multline*}
      yang sama dengan
      $c\cdot f(a_0+a_1x+\dots +a_5x^5)+d\cdot f(b_0+b_1x+\dots +b_5x^5)$.
    \end{answer}
  \item  
    Mengapa dalam \nearbylemma{le:IsoSendsZeroToZero} harus terdapat
    \( \vec{v}\in V \)?
    Dengan kata lain, mengapa $V$ harus tak kosong?
     \begin{answer}
       Tidak ada ruang vektor yang himpunan dasarnya kosong.
       Kita dapat mengambil $\vec{v}$ sebagai vektor nol.
     \end{answer}
  \item  
     Apakah sebarang dua ruang trivial isomorfik?
     \begin{answer}
        Ya; jika kedua ruang tersebut adalah \( \set{\vec{a}} \) dan
        \( \set{\vec{b}} \), pemetaan yang membawa \( \vec{a} \) ke
        \( \vec{b} \) jelas injektif dan surjektif, serta mempertahankan
        sedikit struktur yang ada.
      \end{answer}
  \item 
    Dalam pembuktian \nearbylemma{le:PresStructIffPresCombos}, bagaimana
    dengan kasus tanpa suku (yakni, jika $n$ sama dengan nol)?
    \begin{answer}
       Kombinasi linear dari $n=0$ vektor berjumlah vektor nol, sehingga
       \nearbylemma{le:IsoSendsZeroToZero} menunjukkan bahwa ketiga pernyataan
       tersebut ekuivalen dalam kasus ini.
    \end{answer}
  \item 
    Tunjukkan bahwa setiap isomorfisme \( \map{f}{\polyspace_0}{\Re^1} \)
    berbentuk \( a\mapsto ka \) untuk suatu bilangan riil tak nol \( k \).
    \begin{answer}
      Tinjau basis \( \sequence{1} \) bagi \( \polyspace_0 \), dan misalkan
      \( f(1)\in\Re \) adalah \( k \).
      Untuk sebarang \( a\in\polyspace_0 \), kita memiliki
      \( f(a)=f(a\cdot 1)=af(1)=ak \), sehingga tindakan \( f \) adalah
      perkalian dengan \( k \).
      Perhatikan bahwa \( k\neq 0 \); jika tidak, pemetaan tersebut tidak
      injektif.
      (Selain itu, setiap pemetaan semacam $a\mapsto ka$ merupakan isomorfisme,
      seperti yang mudah diperiksa.)
     \end{answer}
  \item 
    Butir-butir berikut membuktikan bahwa isomorfisme merupakan relasi
    ekuivalensi.
    \begin{exparts}
      \partsitem Tunjukkan bahwa pemetaan identitas
        $\map{\mbox{id}}{V}{V}$ merupakan isomorfisme.
        Jadi, setiap ruang vektor isomorfik dengan dirinya sendiri.
      \partsitem Tunjukkan bahwa jika $\map{f}{V}{W}$ merupakan isomorfisme,
        inversnya $\map{f^{-1}}{W}{V}$ juga merupakan isomorfisme.
        Jadi, jika $V$ isomorfik dengan $W$, maka $W$ juga isomorfik dengan $V$.
      \partsitem Tunjukkan bahwa komposisi isomorfisme merupakan isomorfisme:
        jika $\map{f}{V}{W}$ dan $\map{g}{W}{U}$ merupakan isomorfisme,
        maka $\map{\composed{g}{f}}{V}{U}$ juga merupakan isomorfisme.
        Jadi, jika $V$ isomorfik dengan $W$ dan $W$ isomorfik dengan $U$,
        maka $V$ juga isomorfik dengan $U$.
    \end{exparts}
    \begin{answer}
      Pada setiap butir, dengan mengikuti butir~(2) dari
      \nearbylemma{le:PresStructIffPresCombos}, kita menunjukkan bahwa pemetaan
      mempertahankan struktur dengan menunjukkan bahwa pemetaan tersebut
      mempertahankan kombinasi linear dua unsur domain.
      \begin{exparts}
        \partsitem 
          Pemetaan identitas jelas injektif dan surjektif.
          Pemeriksaan kombinasi linear mudah dilakukan.
          \begin{equation*}
            \mbox{id}(c_1\cdot \vec{v}_1+c_2\cdot \vec{v}_2)
             =c_1\vec{v}_1+c_2\vec{v}_2
             =c_1\cdot \mbox{id}(\vec{v}_1)+c_2\cdot \mbox{id}(\vec{v}_2)
          \end{equation*}
        \partsitem Invers suatu korespondensi juga merupakan korespondensi
          (seperti dinyatakan dalam lampiran), sehingga kita hanya perlu
          memeriksa bahwa invers tersebut mempertahankan kombinasi linear.
          Andaikan \( \vec{w}_1=f(\vec{v}_1) \)
          (sehingga \( f^{-1}(\vec{w}_1)=\vec{v}_1 \)) dan andaikan
          \( \vec{w}_2=f(\vec{v}_2) \).
          \begin{align*}
            f^{-1}(c_1\cdot\vec{w}_1+c_2\cdot\vec{w}_2)
              &=f^{-1}\bigl(\,c_1\cdot f(\vec{v}_1)
                             +c_2\cdot f(\vec{v}_2)\,\bigr) \\
              &=f^{-1}(\,f\bigl(c_1\vec{v}_1+c_2\vec{v}_2)\,\bigr)    \\
              &=c_1\vec{v}_1+c_2\vec{v}_2                             \\ 
              &=c_1\cdot f^{-1}(\vec{w}_1)+c_2\cdot f^{-1}(\vec{w}_2)    
          \end{align*}
        \partsitem Komposisi dua korespondensi merupakan korespondensi
          (seperti dinyatakan dalam lampiran), sehingga kita hanya perlu
          memeriksa bahwa pemetaan komposisinya mempertahankan kombinasi linear.
          \begin{align*}
           \composed{g}{f}\,\bigl(c_1\cdot\vec{v}_1+c_2\cdot\vec{v}_2\bigr)
             &=g\bigl(\,f(c_1\vec{v}_1+c_2\vec{v}_2)\,\bigr)        \\
             &=g\bigl(\,c_1\cdot f(\vec{v}_1)+c_2\cdot f(\vec{v}_2)\,\bigr) \\ 
             &=c_1\cdot g\bigl(f(\vec{v}_1))+c_2\cdot g(f(\vec{v}_2)\bigr)  \\
             &=c_1\cdot\composed{g}{f}\,(\vec{v}_1)
                  +c_2\cdot\composed{g}{f}\,(\vec{v}_2)
          \end{align*}
      \end{exparts}
     \end{answer}
  \item 
    Andaikan \( \map{f}{V}{W} \) mempertahankan struktur.
    Tunjukkan bahwa \( f \) injektif jika dan hanya jika satu-satunya unsur
    \( V \) yang dipetakan oleh \( f \) ke \( \zero_W \) adalah \( \zero_V \).
    \begin{answer}
      Salah satu arah mudah:~menurut definisi, jika \( f \) injektif, maka
      untuk sebarang \( \vec{w}\in W \), paling banyak ada satu
      \( \vec{v}\in V \) yang memenuhi \( f(\vec{v}\,)=\vec{w} \).
      Jadi, khususnya, paling banyak satu unsur \( V \) dipetakan ke
      \( \zero_W \).
      Pembuktian \nearbylemma{le:IsoSendsZeroToZero} tidak menggunakan fakta
      bahwa pemetaan tersebut merupakan korespondensi, sehingga menunjukkan
      bahwa setiap pemetaan \( f \) yang mempertahankan struktur membawa
      \( \zero_V \) ke \( \zero_W \).

      Untuk arah sebaliknya, andaikan satu-satunya unsur \( V \) yang dipetakan
      ke \( \zero_W \) adalah \( \zero_V \).
      Untuk menunjukkan bahwa \( f \) injektif, andaikan
      \( f(\vec{v}_1)=f(\vec{v}_2) \).
      Maka \( f(\vec{v}_1)-f(\vec{v}_2)=\zero_W \), sehingga
      \( f(\vec{v}_1-\vec{v}_2)=\zero_W \).
      Akibatnya \( \vec{v}_1-\vec{v}_2=\zero_V \), sehingga
      $\vec{v}_1=\vec{v}_2$ dan \( f \) injektif.
    \end{answer}
  \item 
    Andaikan \( \map{f}{V}{W} \) merupakan isomorfisme.
    Buktikan bahwa himpunan
    \( \set{\vec{v}_1,\dots,\vec{v}_k}\subseteq V \) bergantung linear
    jika dan hanya jika himpunan bayangan
    \( \set{f(\vec{v}_1),\dots,f(\vec{v}_k)}\subseteq W \)
    bergantung linear.
    \begin{answer}
      Kita akan membuktikan sesuatu yang lebih kuat\Dash bukan hanya keberadaan
      kebergantungan yang dipertahankan oleh isomorfisme, melainkan setiap
      kejadian kebergantungan juga dipertahankan, yaitu,
      \begin{multline*}
        \vec{v}_i=c_1\vec{v}_1+\cdots+c_{i-1}\vec{v}_{i-1}
                   +c_{i+1}\vec{v}_{i+1}+\cdots+c_k\vec{v}_k   \\
        \iff
        f(\vec{v}_i)=c_1f(\vec{v}_1)+\cdots+c_{i-1}f(\vec{v}_{i-1})
                     +c_{i+1}f(\vec{v}_{i+1})+\cdots+c_kf(\vec{v}_k).
      \end{multline*}
      Arah \( \implies \) dari pernyataan ini berlaku berdasarkan butir~(3) dari
      \nearbylemma{le:PresStructIffPresCombos}.
      Arah \( \isimpliedby \) berlaku dengan mengelompokkan kembali
      \begin{align*}
         f(\vec{v}_i)
           &=c_1f(\vec{v}_1)+\dots+c_{i-1}f(\vec{v}_{i-1})
                     +c_{i+1}f(\vec{v}_{i+1})+\dots+c_kf(\vec{v}_k)  \\
           &=f(c_1\vec{v}_1+\dots+c_{i-1}\vec{v}_{i-1}
                     +c_{i+1}\vec{v}_{i+1}+\dots+c_k \vec{v}_k) 
      \end{align*}
      dan menerapkan fakta bahwa $f$ injektif. Jadi, karena kedua vektor
      $\vec{v}_i$ dan
      $c_1\vec{v}_1+\dots+c_{i-1}\vec{v}_{i-1}
              +c_{i+1}\vec{v}_{i+1}+\dots+c_k\vec{v}_k$
      dipetakan oleh $f$ ke bayangan yang sama, keduanya harus sama.
    \end{answer}
  \recommended \item \label{exer:RigidPlaneMapsAutos} 
    Tunjukkan bahwa setiap jenis pemetaan dari
    \nearbyexample{exam:RigidPlaneMapsAutos} merupakan automorfisme.
    \begin{exparts}
      \partsitem Dilatasi $d_s$ oleh skalar tak nol \( s \).
      \partsitem Rotasi $t_\theta$ melalui sudut \( \theta \).
      \partsitem Pencerminan $f_\ell$ terhadap garis yang melalui titik asal.
    \end{exparts}
    \textit{Petunjuk.}
     Untuk butir kedua dan ketiga, koordinat polar berguna.
    \begin{answer}
      \begin{exparts}
        \partsitem 
          Pemetaan ini injektif karena jika
          \( d_s(\vec{v}_1)=d_s(\vec{v}_2) \), maka menurut definisi pemetaan,
          \( s\cdot\vec{v}_1=s\cdot\vec{v}_2 \), sehingga
          \( \vec{v}_1=\vec{v}_2 \), sebab \( s \) tak nol.
          Pemetaan ini surjektif karena sebarang \( \vec{w}\in\Re^2 \)
          merupakan bayangan dari \( \vec{v}=(1/s)\cdot\vec{w} \)
          (sekali lagi, perhatikan bahwa $s$ tak nol).
          (Cara lain untuk melihat bahwa pemetaan ini merupakan korespondensi
          adalah dengan mengamati bahwa pemetaan tersebut memiliki invers:
          invers dari \( d_s \) adalah \( d_{1/s} \).)

          Sebagai penutup, perhatikan bahwa pemetaan ini mempertahankan
          kombinasi linear
          \begin{equation*}
            d_s(c_1\cdot\vec{v}_1+c_2\cdot\vec{v}_2)
            =s(c_1\vec{v}_1+c_2\vec{v}_2)
            =c_1s\vec{v}_1+c_2s\vec{v}_2
            =c_1\cdot d_s(\vec{v}_1)+c_2\cdot d_s(\vec{v}_2)
          \end{equation*}
          sehingga merupakan isomorfisme.
        \partsitem Seperti pada butir sebelumnya, kita dapat menunjukkan bahwa
          pemetaan $t_\theta$ merupakan korespondensi dengan mengamati bahwa
          pemetaan itu memiliki invers, $t_{-\theta}$.

          Secara geometris mudah dilihat bahwa pemetaan ini mempertahankan
          struktur.
          Sebagai contoh, menjumlahkan dua vektor lalu merotasinya memberikan
          hasil yang sama dengan merotasi keduanya terlebih dahulu lalu
          menjumlahkannya.
          Untuk argumen aljabar, tinjau koordinat polar:
          pemetaan \( t_\theta \) membawa vektor dengan titik ujung
          \( (r,\phi) \) ke vektor dengan titik ujung
          \( (r,\phi+\theta) \).
          Selanjutnya, rumus-rumus trigonometri yang dikenal
          $\cos(\phi+\theta)=\cos\phi\,\cos\theta-\sin\phi\,\sin\theta$
          dan
          $\sin(\phi+\theta)=\sin\phi\,\cos\theta+\cos\phi\,\sin\theta$
          menunjukkan cara menyatakan tindakan pemetaan tersebut dalam
          sistem koordinat Kartesius yang biasa.
          \begin{equation*}
            \colvec{x \\ y}=\colvec{r\cos\phi \\ r\sin\phi}
              \mapsunder{t_\theta}
            \colvec{r\cos(\phi+\theta) \\ r\sin(\phi+\theta)}
            =\colvec{x\cos\theta-y\sin\theta \\ x\sin\theta+y\cos\theta}
          \end{equation*}
          Sekarang perhitungan untuk mempertahankan penjumlahan bersifat rutin.
          \begin{multline*}
            \colvec{x_1+x_2 \\ y_1+y_2}
              \mapsunder{t_\theta}
            \colvec{(x_1+x_2)\cos\theta-(y_1+y_2)\sin\theta \\ 
                       (x_1+x_2)\sin\theta+(y_1+y_2)\cos\theta}    \\
            =\colvec{x_1\cos\theta-y_1\sin\theta \\ 
                       x_1\sin\theta+y_1\cos\theta}
            +\colvec{x_2\cos\theta-y_2\sin\theta \\ 
                       x_2\sin\theta+y_2\cos\theta}
          \end{multline*}
          Perhitungan untuk mempertahankan perkalian skalar serupa.
        \partsitem 
          Pemetaan ini merupakan korespondensi karena memiliki invers
          (yaitu dirinya sendiri).

          Seperti pada butir sebelumnya, secara geometris mudah dilihat bahwa
          pemetaan pencerminan mempertahankan struktur:~misalnya, menjumlahkan
          vektor lalu mencerminkannya memberi hasil yang sama dengan
          mencerminkannya terlebih dahulu lalu menjumlahkannya.
          Untuk pembuktian aljabar, andaikan garis $\ell$ memiliki gradien $k$
          (kasus garis dengan gradien tak terdefinisi dapat ditangani sebagai
          kasus terpisah yang mudah).
          Kita dapat mengikuti petunjuk dan menggunakan koordinat polar:~jika
          garis \( \ell \) membentuk sudut \( \phi \) dengan sumbu \( x \),
          tindakan \( f_\ell \) adalah membawa vektor dengan titik ujung
          \( (r\cos\theta,r\sin\theta) \) ke vektor dengan titik ujung
          \( (r\cos(2\phi-\theta),r\sin(2\phi-\theta)) \).
          \begin{center}  \small
            \includegraphics{map/mp/ch3.72}
          \end{center}
          Untuk mengubahnya ke koordinat Kartesius, kita akan menggunakan
          beberapa rumus trigonometri seperti pada butir sebelumnya.
          Pertama, perhatikan bahwa $\cos\phi$ dan $\sin\phi$ dapat ditentukan
          dari gradien $k$ garis tersebut.
          Gambar berikut
          \begin{center}  \small
            \includegraphics{map/mp/ch3.73}
          \end{center}
          memberikan $\cos\phi=1/\sqrt{1+k^2}$ dan
          $\sin\phi=k/\sqrt{1+k^2}$.
          Sekarang,
          \begin{align*}
            \cos(2\phi-\theta)
              &=\cos(2\phi)\,\cos\theta+\sin(2\phi)\,\sin\theta        \\
              &=\left(\cos^2\phi-\sin^2\phi\right)\,\cos\theta
                 +\left(2\sin\phi\cos\phi\right)\,\sin\theta        \\
              &=\left((\frac{1}{\sqrt{1+k^2}})^2
                      -(\frac{k}{\sqrt{1+k^2}})^2\right)
                   \,\cos\theta
                 +\left(2\frac{k}{\sqrt{1+k^2}}\frac{1}{\sqrt{1+k^2}}\right)
                   \,\sin\theta        \\
              &=\left(\frac{1-k^2}{1+k^2}\right)\,\cos\theta
                 +\left(\frac{2k}{1+k^2}\right)\,\sin\theta        
          \end{align*}
          sehingga komponen pertama vektor bayangan adalah sebagai berikut.
          \begin{equation*}
            r\cdot\cos(2\phi-\theta)=\frac{1-k^2}{1+k^2}\cdot x
                                     +\frac{2k}{1+k^2}\cdot y
          \end{equation*}
          Perhitungan serupa menunjukkan bahwa komponen kedua vektor bayangan
          adalah sebagai berikut.
          \begin{equation*}
            r\cdot\sin(2\phi-\theta)=\frac{2k}{1+k^2}\cdot x
                                     -\frac{1-k^2}{1+k^2}\cdot y
          \end{equation*}
          Dengan deskripsi aljabar tentang tindakan $f_\ell$ ini
          \begin{equation*}
            \colvec{x \\ y}
             \mapsunder{f_\ell}
            \colvec{((1-k^2)/(1+k^2))\cdot x
                                     +(2k/(1+k^2))\cdot y \\ 
                   (2k/(1+k^2))\cdot x
                                     -((1-k^2)/(1+k^2))\cdot y}
          \end{equation*}
          pemeriksaan bahwa pemetaan tersebut mempertahankan struktur bersifat
          rutin.
      \end{exparts}  
    \end{answer}
  \item 
    Buat sebuah automorfisme dari \( \polyspace_2 \) selain pemetaan identitas
    dan selain pemetaan pergeseran $p(x)\mapsto p(x-k)$.
    \begin{answer}
       Pertama, pemetaan $p(x)\mapsto p(x+k)$ tidak dihitung karena merupakan
       versi dari $p(x)\mapsto p(x-k)$.
       Berikut sebuah jawaban yang benar (banyak jawaban lain juga benar):
       \( a_0+a_1x+a_2x^2 \mapsto a_2+a_0x+a_1x^2 \).
       Verifikasi bahwa pemetaan ini merupakan isomorfisme berlangsung langsung.
    \end{answer}
  \item
    \begin{exparts}
      \partsitem Tunjukkan bahwa fungsi \( \map{f}{\Re^1}{\Re^1} \)
        merupakan automorfisme jika dan hanya jika berbentuk
        \( x\mapsto kx \) untuk suatu \( k\neq 0 \).
      \partsitem Misalkan \( f \) adalah automorfisme dari \( \Re^1 \)
        sedemikian sehingga \( f(3)=7 \).
        Tentukan \( f(-2) \).
      \partsitem Tunjukkan bahwa fungsi \( \map{f}{\Re^2}{\Re^2} \)
        merupakan automorfisme jika dan hanya jika berbentuk
        \begin{equation*}
          \colvec{x \\ y}
            \mapsto
          \colvec{ax+by \\ cx+dy}
        \end{equation*}
        untuk suatu \( a,b,c,d\in\Re \) dengan \( ad-bc\neq 0 \).
        \textit{Petunjuk.}
        Latihan-latihan dalam subbagian sebelumnya telah menunjukkan bahwa
        \begin{equation*}
          \colvec{a \\ c}\text{ dan }\colvec{b \\ d}
          \text{ bebas linear}
        \end{equation*}
        jika dan hanya jika $ad-bc\neq 0$.
      \partsitem Misalkan \( f \) adalah automorfisme dari \( \Re^2 \) dengan
        \begin{equation*}
          f(\colvec[r]{1 \\ 3})=\colvec[r]{2 \\ -1}
           \quad\text{dan}\quad
          f(\colvec[r]{1 \\ 4})=\colvec[r]{0 \\ 1}.
        \end{equation*}
        Tentukan
        \begin{equation*}
          f(\colvec[r]{0 \\ -1}).
        \end{equation*}
    \end{exparts}
    \begin{answer}
      \begin{exparts}
        \partsitem Untuk arah `hanya jika', misalkan
          \( \map{f}{\Re^1}{\Re^1} \) merupakan isomorfisme.
          Tinjau basis \( \sequence{1}\subseteq\Re^1 \).
          Nyatakan \( f(1) \) dengan \( k \).
          Maka untuk sebarang \( x \), kita memiliki
          \( f(x)=f(x\cdot 1)=x\cdot f(1)=xk \), sehingga tindakan $f$
          adalah perkalian dengan $k$.
          Untuk menyelesaikan arah ini, cukup perhatikan bahwa \( k\neq 0 \);
          jika tidak, $f$ tidak akan injektif.

          Untuk arah `jika', kita hanya perlu memeriksa bahwa pemetaan semacam
          itu merupakan isomorfisme ketika $k\neq 0$.
          Untuk memeriksa sifat injektif, andaikan \( f(x_1)=f(x_2) \),
          sehingga \( kx_1=kx_2 \), lalu bagi dengan faktor tak nol \( k \)
          untuk menyimpulkan $x_1=x_2$.
          Untuk memeriksa sifat surjektif, perhatikan bahwa sebarang
          \( y\in\Re^1 \) merupakan bayangan dari \( x=y/k \)
          (sekali lagi, $k\neq 0$).
          Terakhir, untuk memeriksa bahwa pemetaan semacam itu mempertahankan
          kombinasi linear dua unsur domain, kita memiliki
          \begin{equation*}
            f(c_1x_1+c_2x_2)
            =k(c_1x_1+c_2x_2)
            =c_1kx_1+c_2kx_2
            =c_1f(x_1)+c_2f(x_2)
          \end{equation*}
        \partsitem Berdasarkan butir sebelumnya, tindakan $f$ adalah
          \( x\mapsto (7/3)x \). Jadi \( f(-2)=-14/3 \).
        \partsitem Untuk arah `hanya jika', andaikan
          \( \map{f}{\Re^2}{\Re^2} \) merupakan automorfisme.
          Tinjau basis standar \( \stdbasis_2 \) bagi \( \Re^2 \).
          Misalkan
          \begin{equation*}
            f(\vec{e}_1)=\colvec{a \\ c}
            \quad\text{dan}\quad
            f(\vec{e}_2)=\colvec{b \\ d}.
          \end{equation*}
          Maka tindakan $f$ pada sebarang vektor ditentukan oleh tindakannya
          pada kedua vektor basis.
          \begin{equation*}
            f(\colvec{x \\ y})
            =f(x\cdot\vec{e}_1+y\cdot\vec{e}_2)
            =x\cdot f(\vec{e}_1)+y\cdot f(\vec{e}_2)
            =x\cdot\colvec{a \\ c}+y\cdot\colvec{b \\ d}
            =\colvec{ax+by \\ cx+dy}
          \end{equation*}
          Untuk menyelesaikan arah ini, perhatikan bahwa jika \( ad-bc=0 \),
          yakni jika \( f(\vec{e}_1) \) dan \( f(\vec{e}_2) \) bergantung
          linear, maka \( f \) tidak injektif.

          Untuk arah `jika', kita harus memeriksa bahwa pemetaan tersebut
          merupakan isomorfisme di bawah syarat $ad-bc\neq 0$.
          Pemeriksaan bahwa struktur dipertahankan mudah; di sini kita akan
          menunjukkan bahwa \( f \) merupakan korespondensi.
          Untuk argumen bahwa pemetaan tersebut injektif, andaikan
          \begin{equation*}
            f(\colvec{x_1 \\ y_1})
            =f(\colvec{x_2 \\ y_2})
            \quad\text{sehingga}\quad
            \colvec{ax_1+by_1 \\ cx_1+dy_1}
            =\colvec{ax_2+by_2 \\ cx_2+dy_2}
          \end{equation*}
          Karena \( ad-bc\neq 0 \), sistem yang dihasilkan
          \begin{equation*}
            \begin{linsys}{2}
              a(x_1-x_2)  &+  &b(y_1-y_2)  &=  &0  \\
              c(x_1-x_2)  &+  &d(y_1-y_2)  &=  &0  
            \end{linsys}
          \end{equation*}
          memiliki solusi tunggal, yaitu solusi trivial
          $x_1-x_2=0$ dan $y_1-y_2=0$
          (hal ini mengikuti dari petunjuk).

          Argumen bahwa pemetaan ini surjektif berkaitan erat\Dash sistem
          \begin{equation*}
            \begin{linsys}{2}
              ax_1  &+  &by_1  &=  &x  \\
              cx_1  &+  &dy_1  &=  &y  
            \end{linsys}
          \end{equation*}
          memiliki solusi untuk sebarang \( x \) dan \( y \) jika dan hanya jika
          himpunan
          \begin{equation*}
            \set{\colvec{a \\ c},
                 \colvec{b \\ d} }
          \end{equation*}
          merentang \( \Re^2 \), yakni jika dan hanya jika himpunan ini
          merupakan basis (karena merupakan subhimpunan dua unsur dari
          $\Re^2$), yakni jika dan hanya jika \( ad-bc\neq 0 \).
        \partsitem
          \begin{equation*}
            f(\colvec[r]{0 \\ -1})=f(\colvec[r]{1 \\ 3}-\colvec[r]{1 \\ 4})
            =f(\colvec[r]{1 \\ 3})-f(\colvec[r]{1 \\ 4})
            =\colvec[r]{2 \\ -1}-\colvec[r]{0 \\ 1}
            =\colvec[r]{2 \\ -2}
          \end{equation*}
      \end{exparts}    
    \end{answer}
  \item 
     Lihat \nearbylemma{le:IsoSendsZeroToZero} dan
     \nearbylemma{le:PresStructIffPresCombos}. 
     Temukan dua hal lain yang dipertahankan oleh isomorfisme.
     \begin{answer}
       Terdapat banyak jawaban; dua di antaranya adalah kebebasan linear dan
       subruang.

       Pertama, kita menunjukkan bahwa jika himpunan
       $\set{\vec{v}_1,\dots,\vec{v}_n}$ bebas linear, maka bayangannya
       $\set{f(\vec{v}_1),\dots,f(\vec{v}_n)}$ juga bebas linear.
       Tinjau suatu relasi linear di antara unsur-unsur himpunan bayangan.
       \begin{equation*}
         0=c_1f(\vec{v}_1)+\dots+c_nf(\vec{v}_n)
          =f(c_1\vec{v}_1)+\dots+f(c_n\vec{v}_n)
          =f(c_1\vec{v}_1+\dots+c_n\vec{v}_n)
       \end{equation*}
       Karena pemetaan ini merupakan isomorfisme, pemetaan tersebut injektif.
       Jadi, $f$ hanya memetakan satu vektor dari domain ke vektor nol dalam
       kodomain; artinya, $c_1\vec{v}_1+\dots+c_n\vec{v}_n$ sama dengan vektor nol
       (tentu saja, dalam domain).
       Namun, jika $\set{\vec{v}_1,\dots,\vec{v}_n}$ bebas linear, semua
       koefisien $c$ bernilai nol, sehingga
       $\set{f(\vec{v}_1),\dots,f(\vec{v}_n)}$ juga bebas linear.
       (\textit{Catatan.}
       Ada satu hal kecil mengenai argumen ini yang patut disebutkan.
       Dalam suatu himpunan, pengulangan melebur; artinya, secara ketat
       $\set{\vec{v},\vec{v}}$ merupakan himpunan satu unsur karena kedua hal
       yang dicantumkan di dalamnya sama.
       Namun, perhatikan penggunaan subskrip~$n$ dalam argumen di atas.
       Ketika kita beralih dari himpunan domain
       $\set{\vec{v}_1,\dots,\vec{v}_n}$ ke himpunan bayangan
       $\set{f(\vec{v}_1),\dots,f(\vec{v}_n)}$, tidak terjadi peleburan,
       karena himpunan bayangan tidak memiliki pengulangan, sebab isomorfisme
       $f$ bersifat injektif.)

       Untuk menunjukkan bahwa jika $\map{f}{V}{W}$ merupakan isomorfisme dan
       $U$ adalah subruang dari domain $V$, maka himpunan vektor bayangan
       $f(U)=\set{\vec{w}\in W
                 \suchthat\text{$\vec{w}=f(\vec{u})$ untuk suatu $\vec{u}\in U$}}$
       merupakan subruang dari $W$, kita hanya perlu menunjukkan bahwa
       himpunan tersebut tertutup terhadap kombinasi linear dua unsurnya
       (himpunan ini tak kosong karena memuat bayangan vektor nol).
       Kita memiliki
       \begin{equation*}
         c_1\cdot f(\vec{u}_1)+c_2\cdot f(\vec{u}_2)
        =f(c_1\vec{u}_1)+f(c_2\vec{u}_2)
        =f(c_1\vec{u}_1+c_2\vec{u}_2)
       \end{equation*}
       dan $c_1\vec{u}_1+c_2\vec{u}_2$ merupakan unsur $U$ karena subruang
       tertutup terhadap kombinasi linear.
       Dengan demikian, kombinasi $f(\vec{u}_1)$ dan $f(\vec{u}_2)$ merupakan
       unsur $f(U)$.
      \end{answer}
  \item 
    Kita menunjukkan bahwa isomorfisme dapat disesuaikan: dalam keadaan
    tertentu, jika diberikan vektor-vektor dalam domain dan kodomain, kita dapat
    menghasilkan isomorfisme yang memasangkan vektor-vektor tersebut.
    \begin{exparts}
      \partsitem Misalkan \( B=\sequence{\vec{\beta}_1,\vec{\beta}_2,
                   \vec{\beta}_3} \)
        merupakan basis bagi \( \polyspace_2 \), sehingga sebarang
        \( \vec{p}\in\polyspace_2 \) memiliki representasi tunggal
        sebagai \(
        \vec{p}=c_1\vec{\beta}_1+c_2\vec{\beta}_2+c_3\vec{\beta}_3 \),
        yang kita nyatakan dengan cara berikut.
        \begin{equation*}
          \rep{\vec{p}}{B}=\colvec{c_1 \\ c_2 \\ c_3}
        \end{equation*}
        Tunjukkan bahwa operasi $\rep{\cdot}{B}$ merupakan fungsi dari
        \( \polyspace_2 \) ke \( \Re^3 \)
        (hal ini mencakup menunjukkan bahwa dengan setiap vektor domain
        \( \vec{v}\in\polyspace_2 \)
        terdapat vektor bayangan terkait dalam \( \Re^3 \), dan selanjutnya,
        dengan setiap vektor domain \( \vec{v}\in\polyspace_2 \) terdapat
        paling banyak satu vektor bayangan terkait).
      \partsitem Tunjukkan bahwa fungsi $\rep{\cdot}{B}$ ini injektif dan
        surjektif.
      \partsitem Tunjukkan bahwa fungsi tersebut mempertahankan struktur.
      \partsitem Buat suatu isomorfisme dari \( \polyspace_2 \) ke \( \Re^3 \)
        yang memenuhi spesifikasi berikut.
        \begin{equation*}
          x+x^2\mapsto\colvec[r]{1 \\ 0 \\ 0}
          \quad\text{dan}\quad
          1-x\mapsto\colvec[r]{0 \\ 1 \\ 0}
        \end{equation*}
    \end{exparts}
    \begin{answer}
      \begin{exparts}
        \partsitem Pemasangan
          \begin{equation*}
            \vec{p}=c_1\vec{\beta}_1+c_2\vec{\beta}_2+c_3\vec{\beta}_3
            \mapsunder{\rep{\cdot}{B}}
            \colvec{c_1 \\ c_2 \\ c_3}
          \end{equation*}
          merupakan fungsi jika setiap unsur $\vec{p}$ dari domain dipasangkan
          dengan sekurang-kurangnya satu unsur kodomain, dan jika setiap unsur
          $\vec{p}$ dari domain dipasangkan dengan paling banyak satu unsur
          kodomain.
          Syarat pertama berlaku karena basis $B$ merentang domain\Dash setiap
          $\vec{p}$ dapat ditulis sebagai sekurang-kurangnya satu kombinasi
          linear dari vektor-vektor $\vec{\beta}$.
          Syarat kedua berlaku karena basis $B$ bebas linear\Dash setiap unsur
          $\vec{p}$ dari domain dapat ditulis sebagai paling banyak satu
          kombinasi linear dari vektor-vektor $\vec{\beta}$.
        \partsitem Untuk argumen injektif, jika
          $\rep{\vec{p}}{B}=\rep{\vec{q}}{B}$, yakni jika
          \( \rep{p_1\vec{\beta}_1+p_2\vec{\beta}_2+p_3\vec{\beta}_3}{B}
             =\rep{q_1\vec{\beta}_1+q_2\vec{\beta}_2+q_3\vec{\beta}_3}{B} \)
          maka
          \begin{equation*}
            \colvec{p_1 \\ p_2 \\ p_3}
            =\colvec{q_1 \\ q_2 \\ q_3}
          \end{equation*}
          maka \( p_1=q_1 \), \( p_2=q_2 \), dan \( p_3=q_3 \), yang
          memberikan kesimpulan $\vec{p}=\vec{q}$.
          Oleh karena itu, pemetaan ini injektif.

          Untuk sifat surjektif, kita cukup mencatat bahwa
          \begin{equation*}
            \colvec{a \\ b \\ c}
          \end{equation*}
          sama dengan $\rep{a\vec{\beta}_1+b\vec{\beta}_2+c\vec{\beta}_3}{B}$,
          sehingga sebarang unsur kodomain $\Re^3$ merupakan bayangan suatu
          unsur domain $\polyspace_2$.
        \partsitem Pemetaan ini menghormati penjumlahan dan perkalian skalar
           karena menghormati kombinasi linear dua unsur domain
           (yakni, kita menggunakan butir~(2) dari
           \nearbylemma{le:PresStructIffPresCombos}): jika
           $\vec{p}=p_1\vec{\beta}_1+p_2\vec{\beta}_2+p_3\vec{\beta}_3$ dan
           $\vec{q}=q_1\vec{\beta}_1+q_2\vec{\beta}_2+q_3\vec{\beta}_3$,
           kita memiliki
           \begin{align*}
             \rep{c\cdot\vec{p}+d\cdot\vec{q}}{B}
             &=\rep{\,(cp_1+dq_1)\vec{\beta}_1+(cp_2+dq_2)\vec{\beta}_2
                  +(cp_3+dq_3)\vec{\beta}_3\,}{B}  \\
             &=\colvec{cp_1+dq_1 \\ cp_2+dq_2 \\ cp_3+dq_3}  \\
             &=c\cdot\colvec{p_1 \\ p_2 \\ p_3}  
               +d\cdot\colvec{q_1 \\ q_2 \\ q_3}  \\
             &=c\cdot\rep{\vec{p}}{B}+d\cdot\rep{\vec{q}}{B}
           \end{align*}
        \partsitem Gunakan sebarang basis \( B \) bagi $\polyspace_2$ yang
          dua unsur pertamanya adalah \( x+x^2 \) dan \( 1-x \), misalnya
          \( B=\sequence{x+x^2,1-x,1} \).
      \end{exparts}   
    \end{answer}
  \item 
    Buktikan bahwa suatu ruang berdimensi \( n \) jika dan hanya jika ruang
    tersebut isomorfik dengan \( \Re^n \).
    \textit{Petunjuk.}
    Tetapkan basis $B$ bagi ruang tersebut dan tinjau pemetaan yang membawa
    suatu vektor ke representasinya terhadap $B$.
    \begin{answer}
      Lihat subbagian berikutnya.
    \end{answer}
  \item 
    \textit{(Memerlukan subbagian opsional tentang Mengombinasikan Subruang.)}
    Misalkan \( U \) dan \( W \) adalah ruang-ruang vektor.
    Definisikan ruang vektor baru yang terdiri atas himpunan
    \( U\times W=\set{(\vec{u},\vec{w})  \suchthat  \vec{u}\in U\text{ dan }
                                                   \vec{w}\in W}  \)
    beserta operasi-operasi berikut.
    \begin{equation*}
      (\vec{u}_1,\vec{w}_1)+(\vec{u}_2,\vec{w}_2)=
          (\vec{u}_1+\vec{u}_2,\vec{w}_1+\vec{w}_2)
      \quad\text{dan}\quad
      r\cdot (\vec{u},\vec{w})=(r\vec{u},r\vec{w})
    \end{equation*}
    Ini merupakan ruang vektor, yaitu
    \definend{jumlah langsung eksternal}\index{jumlah langsung!eksternal}%
    \index{jumlah langsung eksternal} dari \( U \) dan \( W \).
    \begin{exparts}
      \partsitem Periksa bahwa himpunan dengan operasi tersebut merupakan
        ruang vektor.
      \partsitem Temukan basis dan dimensi jumlah langsung eksternal
        \( \polyspace_2\times\Re^2 \).
      \partsitem Apa hubungan antara \( \dim(U) \), \( \dim(W) \), dan
        \( \dim(U\times W) \)?
      \partsitem Andaikan \( U \) dan \( W \) merupakan subruang dari ruang
        vektor \( V \) sedemikian sehingga \( V=U\directsum W \)
        (dalam kasus ini kita mengatakan bahwa $V$ merupakan
        \definend{jumlah langsung internal}\index{jumlah langsung!internal}%
        \index{jumlah langsung internal} dari $U$ dan $W$).
        Tunjukkan bahwa pemetaan \( \map{f}{U\times W}{V} \) yang diberikan oleh
        \begin{equation*}
          (\vec{u},\vec{w})\mapsunder{f} \vec{u}+\vec{w}
        \end{equation*}
        merupakan isomorfisme.
        Jadi, jika jumlah langsung internal terdefinisi, maka jumlah langsung
        internal dan eksternal isomorfik.
    \end{exparts}
    \begin{answer}
      \begin{exparts}
        \partsitem Sebagian besar syarat dalam definisi ruang vektor bersifat
          rutin.
          Di sini kita menguraikan verifikasi bagian~(1) dari definisi tersebut.

          Untuk ketertutupan $U\times W$, perhatikan bahwa karena \( U \) dan
          \( W \) tertutup, kita memiliki $\vec{u}_1+\vec{u}_2\in U$ dan
          $\vec{w}_1+\vec{w}_2\in W$, sehingga
          \( (\vec{u}_1+\vec{u}_2,\vec{w}_1+\vec{w}_2)\in U\times W \).
          Komutativitas penjumlahan dalam \( U\times W \) mengikuti dari
          komutativitas penjumlahan dalam \( U \) dan \( W \).
          \begin{equation*}
             (\vec{u}_1,\vec{w}_1)+(\vec{u}_2,\vec{w}_2)=
                 (\vec{u}_1+\vec{u}_2,\vec{w}_1+\vec{w}_2)
                 =(\vec{u}_2+\vec{u}_1,\vec{w}_2+\vec{w}_1)
             =(\vec{u}_2,\vec{w}_2)+(\vec{u}_1,\vec{w}_1)
          \end{equation*}
          Pemeriksaan asosiativitas penjumlahan serupa.
          Unsur nolnya adalah \( (\zero_U,\zero_W)\in U\times W \), dan invers
          aditif dari \( (\vec{u},\vec{w}) \) adalah
          \( (-\vec{u},-\vec{w}) \).

          Pemeriksaan bagian kedua definisi ruang vektor juga langsung.
        \partsitem Berikut adalah suatu basis
          \begin{equation*}
            \sequence{\, (1,\colvec[r]{0 \\ 0}), (x,\colvec[r]{0 \\ 0}),
              (x^2,\colvec[r]{0 \\ 0}), (\zero,\colvec[r]{1 \\ 0}),
              (\zero,\colvec[r]{0 \\ 1}) \,}
          \end{equation*}
          karena terdapat tepat satu cara untuk merepresentasikan sebarang
          unsur $\polyspace_2\times \Re^2$ terhadap himpunan ini;
          berikut sebuah contoh.
          \begin{equation*}
            (3+2x+x^2,\colvec[r]{5 \\ 4})
             =  
              3\cdot(1,\colvec[r]{0 \\ 0})
              +2\cdot(x,\colvec[r]{0 \\ 0})
              +(x^2,\colvec[r]{0 \\ 0}) 
              +5\cdot(\zero,\colvec[r]{1 \\ 0})
              +4\cdot(\zero,\colvec[r]{0 \\ 1}) 
          \end{equation*}
          Dimensi ruang ini adalah lima.
        \partsitem Kita memiliki \( \dim(U\times W)=\dim(U)+\dim(W) \) karena
          berikut merupakan suatu basis.
          \begin{equation*}
            \sequence{ (\vec{\mu}_1,\zero_W), \dots,
              (\vec{\mu}_{\dim(U)},\zero_W), (\zero_U,\vec{\omega}_1),
              \ldots,(\zero_U,\vec{\omega}_{\dim(W)}) }
          \end{equation*}
        \partsitem Kita mengetahui bahwa jika \( V=U\directsum W \), maka setiap
          \( \vec{v}\in V \) dapat ditulis sebagai
          \( \vec{v}=\vec{u}+\vec{w} \) dalam tepat satu cara.
          Inilah yang kita perlukan untuk membuktikan bahwa fungsi yang diberikan
          merupakan isomorfisme.

          Pertama, untuk menunjukkan bahwa \( f \) injektif, kita dapat
          menunjukkan bahwa jika
          \( f\left((\vec{u}_1,\vec{w}_1)\right)
                 =f\left((\vec{u}_2,\vec{w}_2)\right) \), yakni jika
          \( \vec{u}_1+\vec{w}_1=\vec{u}_2+\vec{w}_2 \), maka
          \( \vec{u}_1=\vec{u}_2 \) dan \( \vec{w}_1=\vec{w}_2 \).
          Namun, pernyataan `setiap $\vec{v}$ merupakan jumlah semacam itu
          dalam tepat satu cara' persis merupakan hal yang diperlukan untuk
          menarik kesimpulan ini.
          Demikian pula, argumen bahwa \( f \) surjektif dilengkapi oleh
          pernyataan `setiap $\vec{v}$ merupakan jumlah semacam itu dalam
          sekurang-kurangnya satu cara'.

          Pemetaan ini juga mempertahankan kombinasi linear
          \begin{align*}
            f(\,c_1\cdot(\vec{u}_1,\vec{w}_1)+c_2\cdot (\vec{u}_2,\vec{w}_2)\,)
            &=f(\,(c_1\vec{u}_1+c_2\vec{u}_2,c_1\vec{w}_1+c_2\vec{w}_2)\,) \\
            &=c_1\vec{u}_1+c_2\vec{u}_2+c_1\vec{w}_1+c_2\vec{w}_2    \\
            &=c_1\vec{u}_1+c_1\vec{w}_1+c_2\vec{u}_2+c_2\vec{w}_2    \\
            &=c_1\cdot f(\,(\vec{u}_1,\vec{w}_1)\,)
               +c_2\cdot f(\,(\vec{u}_2,\vec{w}_2)\,)
          \end{align*}
          sehingga merupakan isomorfisme.
      \end{exparts}     
    \end{answer}
\end{exercises}












%\subsection{Real Spaces Represent\index{representative!of isomorphism classes}
%   \protect$n\protect$--Spaces}
\subsection{Dimensi Mencirikan Isomorfisme}

Dalam subbagian sebelumnya, setelah menyatakan definisi isomorfisme,
kita memberikan beberapa hasil yang mendukung gagasan bahwa pemetaan semacam
itu menggambarkan ruang-ruang sebagai ``sama.''
Di sini kita akan mengembangkan intuisi tersebut.
Ketika dua ruang (yang tidak sama) isomorfik, kita memandang keduanya hampir
sama, atau ekuivalen.
Kita akan menyatakan hal itu secara tepat dengan membuktikan bahwa hubungan
`isomorfik dengan' merupakan relasi ekuivalensi.

\begin{lemma}  \label{lem:IsoInvAlsoIso}
%<*lm:IsoInvAlsoIso>
Invers suatu isomorfisme juga merupakan isomorfisme.
%</lm:IsoInvAlsoIso>
\end{lemma}

\begin{proof}
%<*pf:IsoInvAlsoIso0>
Andaikan $V$ isomorfik dengan $W$ melalui $\map{f}{V}{W}$.
Isomorfisme merupakan korespondensi antara himpunan-himpunan, sehingga
$f$ memiliki fungsi invers $\map{f^{-1}}{W}{V}$ yang juga merupakan
korespondensi.\appendrefs{fungsi invers}\spacefactor1000
%</pf:IsoInvAlsoIso0>

%<*pf:IsoInvAlsoIso1>
Kita akan menunjukkan bahwa karena $f$ mempertahankan kombinasi linear,
$f^{-1}$ juga demikian.
Andaikan $\vec{w}_1,\vec{w}_2\in W$.
Karena merupakan isomorfisme, $f$ surjektif, sehingga terdapat
$\vec{v}_1, \vec{v}_2\in V$ sedemikian sehingga
\( \vec{w}_1=f(\vec{v}_1) \) dan \( \vec{w}_2=f(\vec{v}_2) \).
Maka
\begin{multline*}
  f^{-1}(c_1\cdot\vec{w}_1+c_2\cdot\vec{w}_2)
    =f^{-1}\bigl(\,c_1\cdot f(\vec{v}_1)
                             +c_2\cdot f(\vec{v}_2)\,\bigr) \\
    =f^{-1}(\,f\bigl(c_1\vec{v}_1+c_2\vec{v}_2)\,\bigr)    
    =c_1\vec{v}_1+c_2\vec{v}_2                              
    =c_1\cdot f^{-1}(\vec{w}_1)+c_2\cdot f^{-1}(\vec{w}_2)     
\end{multline*}
karena \( f^{-1}(\vec{w}_1)=\vec{v}_1 \) dan
\( f^{-1}(\vec{w}_2)=\vec{v}_2 \).
Dengan demikian, berdasarkan pernyataan kedua
\nearbylemma{le:PresStructIffPresCombos}, pemetaan ini mempertahankan struktur.
%</pf:IsoInvAlsoIso1>
\end{proof}

\begin{theorem} \label{th:IsoEquivRel}
%<*th:IsoEquivRel>
Isomorfisme\index{relasi ekuivalensi!isomorfisme} merupakan relasi ekuivalensi
antara ruang-ruang vektor.
%</th:IsoEquivRel>
\end{theorem}

\begin{proof}
%<*pf:IsoEquivRel0>
Kita harus membuktikan bahwa relasi tersebut simetris, refleksif, dan transitif.

Untuk memeriksa refleksivitas, yaitu bahwa setiap ruang isomorfik dengan
dirinya sendiri, tinjau pemetaan identitas.
Pemetaan tersebut jelas injektif dan surjektif.
Persamaan berikut menunjukkan bahwa pemetaan itu mempertahankan kombinasi linear.
\begin{equation*}
   \mbox{id}(c_1\cdot \vec{v}_1+c_2\cdot \vec{v}_2)
   =c_1\vec{v}_1+c_2\vec{v}_2
   =c_1\cdot \mbox{id}(\vec{v}_1)+c_2\cdot \mbox{id}(\vec{v}_2)
\end{equation*}
%</pf:IsoEquivRel0>

%<*pf:IsoEquivRel1>
Kesimetrian, yakni bahwa jika $V$ isomorfik dengan~$W$, maka $W$ juga
isomorfik dengan~$V$, berlaku berdasarkan \nearbylemma{lem:IsoInvAlsoIso},
karena setiap pemetaan isomorfisme dari $V$ ke $W$ berpasangan dengan
isomorfisme dari $W$ ke $V$.
%</pf:IsoEquivRel1>

%<*pf:IsoEquivRel2>
Sebagai penutup, kita harus memeriksa transitivitas, yakni bahwa jika $V$
isomorfik dengan $W$ dan $W$ isomorfik dengan $U$, maka $V$ isomorfik dengan
$U$.
Misalkan $\map{f}{V}{W}$ dan $\map{g}{W}{U}$ merupakan isomorfisme.
Tinjau komposisinya $\map{\composed{g}{f}}{V}{U}$.
Karena komposisi korespondensi merupakan korespondensi, kita hanya perlu
memeriksa bahwa komposisi tersebut mempertahankan kombinasi linear.
\begin{align*}
  \composed{g}{f}\:\bigl(c_1\cdot\vec{v}_1+c_2\cdot\vec{v}_2\bigr)
     &=g\bigl(\,f(\,c_1\cdot \vec{v}_1+c_2\cdot\vec{v}_2\,)\,\bigr)        \\
     &=g\bigl(\,c_1\cdot f(\vec{v}_1)+c_2\cdot f(\vec{v}_2)\,\bigr) \\ 
     &=c_1\cdot g\bigl(f(\vec{v}_1))+c_2\cdot g(f(\vec{v}_2)\bigr)  \\
     &=c_1\cdot(\composed{g}{f})\,(\vec{v}_1)
                  +c_2\cdot(\composed{g}{f})\,(\vec{v}_2)
\end{align*}
Jadi, komposisi tersebut merupakan isomorfisme.
%</pf:IsoEquivRel2>
\end{proof}

Karena merupakan relasi ekuivalensi, isomorfisme mempartisi semesta ruang
vektor menjadi kelas-kelas:\index{partisi!ke dalam kelas isomorfisme}
setiap ruang berada dalam tepat satu kelas isomorfisme.
\begin{center} \small
  \raisebox{.5in}{\begin{tabular}{l}
                    Semua ruang vektor \\
                    berdimensi hingga:
                  \end{tabular}}
  \includegraphics{map/mp/ch3.17}
  \raisebox{.3in}{\begin{tabular}{l}
                     $V\isomorphicto W$
                  \end{tabular}}
\end{center}
Hasil berikut mencirikan\index{mencirikan}%
\index{isomorfisme!kelas dicirikan oleh dimensi}
kelas-kelas ini berdasarkan dimensi.
Artinya, kita dapat menggambarkan setiap kelas cukup dengan memberikan
bilangan yang merupakan dimensi semua ruang dalam kelas tersebut.
 

\begin{theorem} \label{th:NDimSpaceIsoRN}
%<*th:NDimSpaceIsoRN>
Ruang-ruang vektor isomorfik jika dan hanya jika dimensinya sama.
%</th:NDimSpaceIsoRN>
\end{theorem}

Dalam pernyataan implikasi ganda ini, pembuktian setiap arah melibatkan
gagasan penting, sehingga kita akan membuktikan keduanya secara terpisah.

\begin{lemma}   \label{lem:IsoImpliesSameDim}
%<*lm:IsoImpliesSameDim>
Jika ruang-ruang isomorfik, maka dimensinya sama.
%</lm:IsoImpliesSameDim>
\end{lemma}

\begin{proof}
%<*pf:IsoImpliesSameDim0>
Kita akan menunjukkan bahwa suatu isomorfisme antara dua ruang memberikan
korespondensi antara basis-basisnya.
Artinya, kita akan menunjukkan bahwa jika \( \map{f}{V}{W} \) merupakan
isomorfisme dan basis bagi domain $V$ adalah
\( B=\sequence{\vec{\beta}_1,\dots,\vec{\beta}_n} \), maka bayangannya
\( D=\sequence{f(\vec{\beta}_1),\dots,f(\vec{\beta}_n)} \) merupakan basis
bagi kodomain \( W \).
(Arah lain dari korespondensi tersebut, yakni bahwa untuk sebarang basis dari
$W$, bayangan inversnya merupakan basis bagi $V$, mengikuti fakta bahwa
$f^{-1}$ juga merupakan isomorfisme, sehingga kita dapat menerapkan kalimat
sebelumnya pada $f^{-1}$.)
%</pf:IsoImpliesSameDim0>

%<*pf:IsoImpliesSameDim1>
Untuk melihat bahwa \( D \) merentang \( W \), tetapkan sebarang
\( \vec{w}\in W \).
Karena \( f \) merupakan isomorfisme, pemetaan tersebut surjektif, sehingga
terdapat \( \vec{v}\in V \) dengan \( \vec{w}=f(\vec{v}) \).
Jabarkan \( \vec{v} \) sebagai kombinasi vektor-vektor basis.
\begin{equation*}
  \vec{w}=f(\vec{v})
  =f(v_1\vec{\beta}_1+\dots+v_n\vec{\beta}_n)  
  =v_1\cdot f(\vec{\beta}_1)+\dots+v_n\cdot f(\vec{\beta}_n)
\end{equation*}
Untuk kebebasan linear $D$, jika
\begin{equation*}
  \zero_W
  =c_1f(\vec{\beta}_1)+\dots+c_nf(\vec{\beta}_n)  
  =f(c_1\vec{\beta}_1+\dots+c_n\vec{\beta}_n)
\end{equation*}
maka, karena \( f \) injektif sehingga satu-satunya vektor yang dibawa ke
\( \zero_W \) adalah \( \zero_V \), kita memiliki
\( \zero_V=c_1\vec{\beta}_1+\dots+c_n\vec{\beta}_n \), yang menyiratkan
bahwa semua koefisien \( c \) bernilai nol.
%</pf:IsoImpliesSameDim1>
\end{proof}

\begin{lemma} \label{lem:EqDimImpIso}
%<*lm:EqDimImpIso>
Jika ruang-ruang memiliki dimensi yang sama, maka ruang-ruang tersebut isomorfik.
%</lm:EqDimImpIso>
\end{lemma}

\begin{proof}
%<*pf:EqDimImpIso0>
Kita akan membuktikan bahwa setiap ruang berdimensi~$n$ isomorfik dengan
$\Re^n$.
Selanjutnya, berdasarkan transitivitas yang ditunjukkan dalam
\nearbytheorem{th:IsoEquivRel}, semua ruang semacam itu isomorfik satu sama lain.
%</pf:EqDimImpIso0>

%<*pf:EqDimImpIso1>
Misalkan $V$ berdimensi \( n \).
Tetapkan basis \( B=\sequence{\vec{\beta}_1,\dots,\vec{\beta}_n} \) bagi
domain $V$.
Tinjau operasi merepresentasikan unsur-unsur $V$ terhadap $B$ sebagai fungsi
dari $V$ ke $\Re^n$.
\begin{equation*}
   \vec{v}=v_1\vec{\beta}_1+\dots+v_n\vec{\beta}_n
   \,\mapsunder{\text{\small Rep}_B}\,\colvec{v_1 \\ \vdots \\ v_n}
\end{equation*}
%</pf:EqDimImpIso1>
Fungsi ini \definend{terdefinisi dengan baik}%
\appendrefs{terdefinisi dengan baik}\spacefactor1000 %
\index{terdefinisi dengan baik}
karena setiap \( \vec{v} \) memiliki tepat satu representasi semacam itu
(lihat \nearbyremark{not:WellDefFcns} setelah pembuktian ini).

%<*pf:EqDimImpIso2>
Fungsi ini injektif karena jika
\begin{equation*}
   \text{Rep}_B(u_1\vec{\beta}_1+\dots+u_n\vec{\beta}_n)
     =\text{Rep}_B(v_1\vec{\beta}_1+\dots+v_n\vec{\beta}_n)
\end{equation*}
maka
\begin{equation*}
  \colvec{u_1 \\ \vdots \\ u_n}
    =
  \colvec{v_1 \\ \vdots \\ v_n}
\end{equation*}
sehingga \( u_1=v_1 \), \ldots, $u_n=v_n$, yang menyiratkan bahwa argumen
semula \( u_1\vec{\beta}_1+\dots+u_n\vec{\beta}_n \) dan
\( v_1\vec{\beta}_1+\dots+v_n\vec{\beta}_n\) sama.
%</pf:EqDimImpIso2>

%<*pf:EqDimImpIso3>
Fungsi ini surjektif; sebarang unsur $\Re^n$
\begin{equation*}
   \vec{w}=\colvec{w_1 \\ \vdots \\ w_n}
\end{equation*}
merupakan bayangan suatu \( \vec{v}\in V \), yaitu
\( \vec{w}=\rep{w_1\vec{\beta}_1+\dots+w_n\vec{\beta}_n}{B}  \).
%</pf:EqDimImpIso3>

%<*pf:EqDimImpIso4>
Terakhir, fungsi ini mempertahankan struktur.
\begin{align*}
  \rep{r\cdot\vec{u}+s\cdot\vec{v}}{B}
  &=\rep{\,(ru_1+sv_1)\vec{\beta}_1+\dots+(ru_n+sv_n)\vec{\beta}_n\,}{B}  \\
  &=\colvec{ru_1+sv_1 \\ \vdots \\ ru_n+sv_n}  \\
  &=r\cdot\colvec{u_1 \\ \vdots \\ u_n}+s\cdot\colvec{v_1 \\ \vdots \\ v_n} \\
  &=r\cdot\rep{\vec{u}}{B}+s\cdot\rep{\vec{v}}{B}
\end{align*}

Oleh karena itu, $\mbox{Rep}_B$ merupakan isomorfisme.
Akibatnya, setiap ruang berdimensi $n$ isomorfik dengan $\Re^n$.
%</pf:EqDimImpIso4>
\end{proof}

\begin{remark}
Ketika memperkenalkan notasi $\mbox{Rep}_B$ bagi vektor pada
halaman~\pageref{rem:RepNotation}, kita mencatat bahwa notasi tersebut tidak
standar dan mengatakan bahwa salah satu keunggulannya adalah lebih sulit
terlewatkan.
Di sini kita melihat keunggulan lainnya:~notasi ini menyatakan secara eksplisit
bahwa $\mbox{Rep}_B$ merupakan fungsi dari~$V$ ke $\R^n$.
\end{remark}

\begin{remark} \label{not:WellDefFcns}
\index{terdefinisi dengan baik}
Pembuktian tersebut memuat sebuah kalimat tentang `terdefinisi dengan baik.'
Maksudnya adalah bahwa untuk menjadi isomorfisme, $\mbox{Rep}_B$ harus
merupakan fungsi.
Definisi fungsi mensyaratkan bahwa bagi setiap masukan, keluaran yang terkait
harus ada dan harus ditentukan oleh masukannya.
Jadi, kita harus memeriksa bahwa setiap $\vec{v}$ dipasangkan dengan
sekurang-kurangnya satu~$\mbox{Rep}_B(\vec{v})$, dan tidak lebih dari satu.

Dalam pembuktian, kita menyatakan unsur-unsur $\vec{v}$ dari ruang domain
sebagai kombinasi unsur-unsur basis~$B$, lalu memasangkan $\vec{v}$ dengan
vektor kolom koefisien.
Keberadaan sekurang-kurangnya satu penjabaran bagi setiap~$\vec{v}$ berlaku
karena $B$ merupakan basis sehingga merentang ruang tersebut.

Persoalan bahwa tidak boleh ada lebih dari satu unsur kodomain yang terkait
lebih halus.
Contoh pembanding ketika sebuah pemasangan gagal memenuhi syarat keluaran
tunggal ini akan memperjelas persoalan tersebut.
Misalkan domainnya \( \polyspace_2 \), dan tinjau suatu himpunan yang bukan
basis (himpunan itu tidak bebas linear, meskipun merentang ruang tersebut).
\begin{equation*}
 A=\set{1+0x+0x^2,
              0+1x+0x^2,
              0+0x+1x^2,
              1+1x+2x^2}
\end{equation*}
Namakan polinomial-polinomial tersebut $\vec{\alpha}_1$, \ldots,
$\vec{\alpha}_4$.
Berbeda dengan keadaan ketika himpunan itu merupakan basis, di sini dapat
terdapat lebih dari satu pernyataan vektor domain dalam bentuk
% \( \vec{v}=c_1\vec{\alpha}_1+c_2\vec{\alpha}_2+
%              c_3\vec{\alpha}_3+c_4\vec{\alpha}_4 \)
unsur-unsur himpunan tersebut.
Sebagai contoh, tinjau \( \vec{v}=1+x+x^2 \).
Berikut dua penjabaran yang berbeda.
\begin{equation*}
  \vec{v}=1\vec{\alpha}_1+1\vec{\alpha}_2+1\vec{\alpha}_3+0\vec{\alpha}_4 
  \qquad
  \vec{v}=0\vec{\alpha}_1+0\vec{\alpha}_2-
                1\vec{\alpha}_3+1\vec{\alpha}_4 
\end{equation*}
Jadi, vektor masukan~$\vec{v}$ ini dipasangkan dengan lebih dari satu kolom.
\begin{equation*}
  \colvec[r]{1 \\ 1 \\ 1 \\ 0}
  \qquad
  \colvec[r]{0 \\ 0 \\ -1 \\ 1}
\end{equation*}
Dengan demikian, dengan~$A$ pemasangan tersebut tidak terdefinisi dengan baik.
(Persoalannya adalah bahwa $A$ tidak bebas linear; untuk menunjukkan
ketunggalan, pembuktian Teorema~Dua.III.\ref{th:BasisIffUniqueRepWRT}
hanya menggunakan kebebasan linear.)

Secara umum, setiap kali kita mendefinisikan fungsi, kita harus memeriksa bahwa
nilai-nilai keluarannya terdefinisi dengan baik.
Biasanya syarat ini sangat jelas, tetapi dalam pembuktian di atas syarat
tersebut perlu diverifikasi.
Lihat \nearbyexercise{exer:FcnWellDef}.
\end{remark}

\begin{corollary} \label{co:FiniteDimensionalIsoToReN}
%<*co:FiniteDimensionalIsoToReN>
Setiap ruang vektor berdimensi hingga isomorfik dengan tepat satu dari
ruang-ruang $\Re^n$.
%</co:FiniteDimensionalIsoToReN>
\end{corollary}

Hal ini memberi kita kumpulan perwakilan kelas-kelas isomorfisme.
\begin{center}
  \raisebox{.5in}{\begin{tabular}{l}
                    Semua ruang vektor \\
                    berdimensi hingga:
                  \end{tabular}}
  \includegraphics{map/mp/ch3.18}
  \raisebox{.3in}{\begin{tabular}{l}
                    Satu perwakilan \\ per kelas
                  \end{tabular}}
\end{center}

Pembuktian-pembuktian di atas memadatkan banyak gagasan dalam ruang yang kecil.
Sepanjang sisa bab ini, kita akan meninjau kembali dan mengembangkan
gagasan-gagasan tersebut.
Sebagai gambaran awal, di sini kita akan memperluas pembuktian
\nearbylemma{lem:EqDimImpIso}.

\begin{example}  \label{ex:ExtendLinearlyMatrixMap}
Ruang $\matspace_{\nbyn{2}}$ yang terdiri atas matriks $\nbyn{2}$ isomorfik
dengan $\Re^4$.
Dengan basis domain berikut
\begin{equation*}
  B=\sequence{\begin{mat}[r]
                1  &0  \\
                0  &0
              \end{mat},
              \begin{mat}[r]
                0  &1  \\
                0  &0
              \end{mat},
              \begin{mat}[r]
                0  &0  \\
                1  &0
              \end{mat},
              \begin{mat}[r]
                0  &0  \\
                0  &1
              \end{mat} }
\end{equation*}
isomorfisme yang diberikan dalam lema, yaitu pemetaan representasi
$f_1=\mbox{Rep}_B$, memindahkan entri-entri tersebut.
\begin{equation*}
  \begin{mat}
    a  &b  \\
    c  &d
  \end{mat}
  \mapsunder{f_1}
  \colvec{a \\ b \\ c \\ d}
\end{equation*}
Salah satu cara memahami pemetaan $f_1$ adalah sebagai berikut:~tetapkan basis
$B$ bagi domain, gunakan basis standar $\stdbasis_4$ bagi kodomain, lalu
pasangkan $\vec{\beta}_1$ dengan $\vec{e}_1$, $\vec{\beta}_2$ dengan
$\vec{e}_2$, dan seterusnya.
Kemudian perluas pemasangan ini ke semua unsur kedua ruang.
\begin{equation*}
  \begin{mat}
    a  &b  \\
    c  &d
  \end{mat}
  =a\vec{\beta}_1+b\vec{\beta}_2+c\vec{\beta}_3+d\vec{\beta}_4
  \;\;\mapsunder{f_1}\;\;
  a\vec{e}_1+b\vec{e}_2+c\vec{e}_3+d\vec{e}_4
  =\colvec{a \\ b \\ c \\ d}
\end{equation*}

Kita dapat melakukan hal yang sama dengan basis-basis berbeda, misalnya dengan
mengambil basis berikut bagi domain.
\begin{equation*}
  A=\sequence{
              \begin{mat}[r]
                2  &0  \\
                0  &0
              \end{mat},
              \begin{mat}[r]
                0  &2  \\
                0  &0
              \end{mat},
              \begin{mat}[r]
                0  &0  \\
                2  &0
              \end{mat},
              \begin{mat}[r]
                0  &0  \\
                0  &2
              \end{mat} }
\end{equation*}
Memasangkan unsur-unsur yang berkorespondensi dari $A$ dan $\stdbasis_4$
memberikan
\begin{multline*}
  \begin{mat}
    a  &b  \\
    c  &d
  \end{mat}
  =(a/2)\vec{\alpha}_1+(b/2)\vec{\alpha}_2
      +(c/2)\vec{\alpha}_3+(d/2)\vec{\alpha}_4          \\
  \mapsunder{f_2}\;\;
  (a/2)\vec{e}_1+(b/2)\vec{e}_2+(c/2)\vec{e}_3+(d/2)\vec{e}_4
  =\colvec{a/2 \\ b/2 \\ c/2 \\ d/2}
\end{multline*}
menghasilkan isomorfisme yang berbeda dari $f_1$.

Pemetaan sebelumnya muncul dengan mengubah basis domain.
Kita juga dapat mengubah basis kodomain.
Kembali ke basis $B$ di atas dan gunakan basis berikut bagi kodomain.
\begin{equation*}
  D=\sequence{\colvec[r]{1 \\ 0 \\ 0 \\ 0},
                \colvec[r]{0 \\ 1 \\ 0 \\ 0},
                \colvec[r]{0 \\ 0 \\ 0 \\ 1},
                \colvec[r]{0 \\ 0 \\ 1 \\ 0}}
\end{equation*}
Pasangkan $\vec{\beta}_1$ dengan $\vec{\delta}_1$, dan seterusnya.
Memperluas pemasangan tersebut memberikan isomorfisme lain.
\begin{equation*}
  \begin{mat}
     a  &b  \\
     c  &d  
   \end{mat}
  =a\vec{\beta}_1+b\vec{\beta}_2+c\vec{\beta}_3+d\vec{\beta}_4
  \;\;\mapsunder{f_3}\;\;
  a\vec{\delta}_1+b\vec{\delta}_2+c\vec{\delta}_3+d\vec{\delta}_4
  =\colvec{a \\ b \\ d \\ c}
\end{equation*}
\end{example}

Kita menutup dengan sebuah rangkuman.
Ingat bahwa bab pertama mendefinisikan dua matriks sebagai ekuivalen baris
jika keduanya dapat diperoleh satu sama lain melalui operasi baris.
Di sana kita menunjukkan bahwa relasi tersebut merupakan relasi ekuivalensi,
sehingga kumpulan matriks dipartisi ke dalam kelas-kelas; semua matriks yang
saling ekuivalen baris berada dalam satu kelas.
Selanjutnya, untuk memahami matriks mana yang berada dalam setiap kelas,
kita memberikan perwakilan kelas-kelas tersebut, yaitu matriks bentuk eselon
tereduksi.

Dalam bagian ini kita mengikuti pola tersebut, kecuali bahwa gagasan ``sama''
di sini adalah isomorfisme ruang vektor.
Kita mendefinisikannya dan menetapkan beberapa sifat, termasuk bahwa
isomorfisme merupakan relasi ekuivalensi.
Kemudian, seperti sebelumnya, kita menyusun daftar perwakilan kelas untuk
membantu memahami partisi tersebut\Dash partisi itu mengklasifikasikan ruang
vektor berdasarkan dimensi.

Dalam Bab Dua, melalui definisi ruang vektor, kita tampaknya telah membuka
kajian terhadap banyak contoh struktur baru selain ruang-ruang $\Re^n$ yang
sudah dikenal.
Sekarang kita mengetahui bahwa kenyataannya tidak demikian.
Setiap ruang vektor berdimensi hingga sesungguhnya ``sama'' dengan suatu
ruang riil.




\begin{exercises}
  \recommended \item 
    Tentukan apakah ruang-ruang berikut isomorfik.
    \begin{exparts*}
       \partsitem \( \Re^2 \),~\( \Re^4 \)
       \partsitem \( \polyspace_5 \),~\( \Re^5 \)
       \partsitem \( \matspace_{\nbym{2}{3}} \),~\( \Re^6 \)
       \partsitem \( \polyspace_5 \),~\( \matspace_{\nbym{2}{3}} \)
       \partsitem \( \matspace_{\nbym{2}{k}} \),~\( \matspace_{\nbym{k}{2}} \)
    \end{exparts*}
    \begin{answer}
       Setiap pasangan ruang isomorfik jika dan hanya jika keduanya memiliki
       dimensi yang sama.
       Jika terdapat isomorfisme, kita dapat menyatakan suatu pemetaan, tetapi
       hal itu tidak mutlak diperlukan.
       \begin{exparts}
         \partsitem Tidak, dimensinya berbeda.
         \partsitem Tidak, dimensinya berbeda.
         \partsitem Ya, dimensinya sama.
           Berikut salah satu isomorfismenya.
           \begin{equation*}
             \begin{mat}
               a  &b  &c  \\
               d  &e  &f
             \end{mat}
             \mapsto
             \colvec{a \\ \vdots \\ f}
           \end{equation*}
         \partsitem Ya, dimensinya sama.
           Berikut suatu isomorfisme.
           \begin{equation*}
             a+bx+\cdots+fx^5
             \mapsto
             \begin{mat}
               a  &b  &c  \\
               d  &e  &f
             \end{mat}
           \end{equation*}
         \partsitem Ya, keduanya berdimensi \( 2k \).
      \end{exparts}    
    \end{answer}
  \item Ruang-ruang manakah berikut ini yang saling isomorfik?
    \begin{exparts*}
      \partsitem $\Re^3$
      \partsitem $\matspace_{\nbyn{2}}$
      \partsitem $\polyspace_3$
      \partsitem $\Re^4$
      \partsitem $\polyspace_2$
    \end{exparts*}
    \begin{answer}
      Cukup tentukan dimensi setiap ruang, misalnya dengan menemukan suatu
      basis; ruang-ruang dengan dimensi yang sama kemudian isomorfik.
      Berikut daftar dimensi setiap ruang.
      \begin{exparts*}
        \partsitem $3$
        \partsitem $4$
        \partsitem $4$
        \partsitem $4$
        \partsitem $3$
      \end{exparts*}
    \end{answer}
  \recommended \item 
    Tinjau isomorfisme
    \( \map{\rep{\cdot}{B}}{\polyspace_1}{\Re^2} \) 
    dengan \( B=\sequence{1,1+x} \).
    Tentukan bayangan setiap unsur domain berikut.
    \begin{exparts*}
      \partsitem \( 3-2x \);
      \partsitem \( 2+2x \);
      \partsitem \( x \)
    \end{exparts*}
    \begin{answer}
      \begin{exparts*}
        \partsitem \( \rep{3-2x}{B}=\colvec[r]{5 \\ -2} \)
        \partsitem \( \colvec[r]{0 \\ 2} \)
        \partsitem \( \colvec[r]{-1 \\ 1} \)
      \end{exparts*}  
     \end{answer}
  \item Untuk nilai $n$ berapakah ruang tersebut isomorfik dengan~$\Re^n$?
    \begin{exparts}
      \partsitem $\polyspace_4$
      \partsitem $\polyspace_1$
      \partsitem $\matspace_{\nbym{2}{3}}$
      \partsitem bidang $2x-y+z=0$ yang merupakan subhimpunan dari~$\Re^3$
      \partsitem ruang vektor kombinasi linear tiga huruf
         $\set{ax+by+cz\suchthat a,b,c\in\Re}$
    \end{exparts}
    \begin{answer}
      Untuk setiap ruang, cara termudah ialah menentukan dimensinya dengan
      menemukan suatu basis.
      Untuk setiap basis yang diberikan di bawah ini, kita mengabaikan
      verifikasi bahwa himpunan tersebut benar-benar merupakan basis.
      \begin{exparts}
        \partsitem Ruang ini isomorfik dengan $\Re^5$.
          Salah satu basis bagi $\polyspace_4$ adalah
          $\set{x^4,x^3,x^2,x,1}$, sehingga ruang tersebut berdimensi~$5$.
        \partsitem Ruang ini isomorfik dengan $\Re^2$ karena salah satu basis
          bagi ruang $\polyspace_1=\set{a+bx\suchthat a,b\in\Re}$ adalah
          $\set{1,x}$.
        \partsitem Ruang ini isomorfik dengan $\Re^6$.
          Salah satu basisnya terdiri atas enam matriks berikut.
          \begin{equation*}
            \begin{mat}
              1 &0 &0 \\
              0 &0 &0
            \end{mat},\; 
            \begin{mat}
              0 &1 &0 \\
              0 &0 &0
            \end{mat},\;
            \cdots\;
            \begin{mat}
              0 &0 &0 \\
              0 &0 &1
            \end{mat}
          \end{equation*}
        \partsitem Ruang ini merupakan bidang, sehingga isomorfik
          dengan~$\Re^2$.
          Untuk jawaban yang lebih lengkap, memparameterkan bidang tersebut
          memberikan deskripsi vektor berikut
          \begin{equation*}
            \set{\colvec{x \\ y \\ z}=\colvec{1/2 \\ 1 \\ 0}y
                                      +\colvec{-1/2 \\ 0 \\ 1}z
                   \suchthat  y,z\in\Re}
          \end{equation*}
          sehingga ruang itu memiliki basis yang terdiri atas kedua vektor
          tersebut.
        \partsitem Ruang ini isomorfik dengan $\Re^3$.
          Salah satu basisnya adalah himpunan kombinasi linear
          $\set{x,y,z}$, yaitu $\set{x+0y+0z, 0x+y+0z,0x+0y+z}$.
      \end{exparts}
    \end{answer}
  \recommended \item
    Tunjukkan bahwa jika \( m\neq n \), maka
    \( \Re^m\not\isomorphicto\Re^n \).
     \begin{answer}
       Keduanya memiliki dimensi yang berbeda.
     \end{answer}
  \recommended \item
    Apakah \( \matspace_{\nbym{m}{n}}\isomorphicto\matspace_{\nbym{n}{m}} \)?
    \begin{answer}
      Ya, keduanya berdimensi \( mn \).
    \end{answer}
  \recommended \item
    Apakah sebarang dua bidang yang melalui titik asal dalam \( \Re^3 \)
    isomorfik?
    \begin{answer}
      Ya, sebarang dua bidang (tak degenerat) sama-sama merupakan ruang vektor
      berdimensi dua.
    \end{answer}
  \item 
     Temukan himpunan perwakilan kelas ekuivalensi selain himpunan
     ruang-ruang \( \Re^n \).
     \begin{answer}
        Terdapat banyak jawaban; salah satunya adalah himpunan ruang-ruang
        \( \polyspace_k \) (dengan mengambil \( \polyspace_{-1} \) sebagai
        ruang vektor trivial).
     \end{answer}
  \item 
    Benar atau salah:~antara sebarang ruang berdimensi \( n \) dan \( \Re^n \)
    terdapat tepat satu isomorfisme.
    \begin{answer}
      Salah (kecuali ketika \( n=0 \)).
      Sebagai contoh, jika \( \map{f}{V}{\Re^n} \) merupakan isomorfisme,
      mengalikannya dengan sebarang skalar tak nol selain \(1\) memberikan
      isomorfisme lain yang berbeda.
      (Antara ruang-ruang trivial, isomorfismenya tunggal; satu-satunya
      pemetaan yang mungkin adalah $\zero_V\mapsto \zero_W$.)
    \end{answer}
  \item 
    Dapatkah suatu ruang vektor isomorfik dengan salah satu subruang sejatinya?
    \begin{answer}
      Dapat, jika ruangnya berdimensi tak hingga. Sebagai contoh, pemetaan
      \( \map{T}{\polyspace}{x\polyspace} \) yang diberikan oleh
      \( T(p(x))=xp(x) \) merupakan isomorfisme dari ruang semua polinomial
      ke subruang sejatinya \(x\polyspace\).
      Namun, dalam konteks berdimensi hingga buku ini, hal tersebut tidak
      mungkin: subruang sejati memiliki dimensi yang lebih kecil secara ketat
      daripada ruang induknya.
    \end{answer}
  \recommended \item 
    Subbagian ini menunjukkan bahwa bagi setiap isomorfisme, pemetaan inversnya
    juga merupakan isomorfisme.
    Subbagian ini juga menunjukkan bahwa untuk basis tetap \( B \) dari ruang
    vektor berdimensi \( n \), \( V \), pemetaan
    \( \map{\text{Rep}_B}{V}{\Re^n} \) merupakan isomorfisme.
    Tentukan invers pemetaan ini.
    \begin{answer}
      Jika \( B=\sequence{\vec{\beta}_1,\ldots,\vec{\beta}_n} \), inversnya
      adalah sebagai berikut.
      \begin{equation*}
        \colvec{c_1 \\ \vdots \\ c_n}
        \mapsto c_1\vec{\beta}_1+\cdots+c_n\vec{\beta}_n
      \end{equation*}  
    \end{answer}
  \recommended \item 
    Buktikan fakta-fakta berikut tentang matriks.
    \begin{exparts}
      \partsitem Ruang baris suatu matriks isomorfik dengan ruang kolom dari
        transposnya.
      \partsitem Ruang baris suatu matriks isomorfik dengan ruang kolomnya.
    \end{exparts}
    \begin{answer}
      Ketiga ruang tersebut memiliki dimensi yang sama dengan rank matriks.
    \end{answer}
  \item \label{exer:FcnWellDef}
    Tunjukkan bahwa fungsi dari \nearbytheorem{th:NDimSpaceIsoRN}
    terdefinisi dengan baik.
     \begin{answer}
        Kita harus menunjukkan bahwa jika \( \vec{a}=\vec{b} \), maka
        \( f(\vec{a})=f(\vec{b}) \).
        Maka andaikan
        $a_1\vec{\beta}_1+\dots+a_n\vec{\beta}_n
           =b_1\vec{\beta}_1+\dots+b_n\vec{\beta}_n$.
        Setiap vektor dalam ruang vektor (di sini, ruang domain) memiliki
        representasi tunggal sebagai kombinasi linear vektor-vektor basis,
        sehingga kita dapat menyimpulkan bahwa \( a_1=b_1 \), \ldots,
        \(a_n=b_n \).
        Jadi,
        \begin{equation*}
          f(\vec{a})
          =\colvec{a_1 \\ \vdots \\ a_n}=\colvec{b_1 \\ \vdots \\ b_n}
          =f(\vec{b})          
        \end{equation*}
        sehingga fungsi tersebut terdefinisi dengan baik.
     \end{answer}
  \item 
     Apakah pembuktian \nearbytheorem{th:NDimSpaceIsoRN} sah ketika \( n=0 \)?
     \begin{answer}
       Ya, karena ruang berdimensi nol merupakan ruang trivial.
     \end{answer}
  \item 
    Untuk setiap himpunan berikut, tentukan apakah himpunan tersebut merupakan
    himpunan perwakilan kelas isomorfisme.
    \begin{exparts}
      \partsitem \( \set{\C^k \suchthat k\in\N} \)
      \partsitem $\set{\polyspace_k\suchthat k\in \set{-1,0,1,\ldots}}$
      \partsitem \( \set{\matspace_{\nbym{m}{n}}\suchthat 
                   m,n\in\N} \)
    \end{exparts}
    \begin{answer}
      \begin{exparts}
        \partsitem Tidak, kumpulan ini tidak memiliki ruang berdimensi ganjil.
        \partsitem Ya, karena $\polyspace_{k}\isomorphicto\Re^{k+1}$.
        \partsitem Tidak, sebagai contoh,
          \( \matspace_{\nbym{2}{3}}\isomorphicto\matspace_{\nbym{3}{2}} \).
      \end{exparts}  
    \end{answer}
  \item
    Misalkan \( f \) merupakan korespondensi antara ruang-ruang vektor
    \( V \) dan \( W \) (yakni, pemetaan yang injektif dan surjektif).
    Tunjukkan bahwa ruang \( V \) dan \( W \) isomorfik melalui \( f \)
    jika dan hanya jika terdapat basis \( B\subset V \) dan \( D\subset W \)
    sedemikian sehingga vektor-vektor yang berkorespondensi memiliki koordinat
    yang sama:
    \( \rep{\vec{v}}{B}=\rep{f(\vec{v})}{D} \).
    \begin{answer}
       Salah satu arah mudah:~jika kedua ruang isomorfik melalui \( f \), maka
       untuk sebarang basis \( B\subseteq V \), himpunan \( D=f(B) \) juga
       merupakan basis (hal ini ditunjukkan dalam
       \nearbylemma{lem:IsoImpliesSameDim}).
       Pemeriksaan bahwa vektor-vektor yang berkorespondensi memiliki koordinat
       yang sama:
       \( f(c_1\vec{\beta}_1+\dots+c_n\vec{\beta}_n)
          =c_1f(\vec{\beta}_1)+\dots+c_nf(\vec{\beta}_n)
          =c_1\vec{\delta}_1+\dots+c_n\vec{\delta}_n   \)
       bersifat rutin.

       Untuk arah lainnya, andaikan terdapat basis-basis sedemikian sehingga
       vektor-vektor yang berkorespondensi memiliki koordinat yang sama terhadap
       basis-basis tersebut.
       Karena \( f \) merupakan korespondensi, untuk menunjukkan bahwa
       pemetaan itu merupakan isomorfisme, kita hanya perlu menunjukkan bahwa
       pemetaan tersebut mempertahankan struktur.
       Karena \( \rep{\vec{v}\,}{B}=\rep{f(\vec{v}\,)}{D} \), pemetaan \( f \)
       mempertahankan struktur jika dan hanya jika representasi mempertahankan
       penjumlahan:
       \( \rep{\vec{v}_1+\vec{v}_2}{B}=\rep{\vec{v}_1}{B}+\rep{\vec{v}_2}{B} \)
       dan perkalian skalar:
       \( \rep{r\cdot\vec{v}\,}{B}=r\cdot\rep{\vec{v}\,}{B} \)
       Perhitungan penjumlahannya adalah sebagai berikut:
       \( (c_1+d_1)\vec{\beta}_1+\dots+(c_n+d_n)\vec{\beta}_n
          =c_1\vec{\beta}_1+\dots+c_n\vec{\beta}_n
          +d_1\vec{\beta}_1+\dots+d_n\vec{\beta}_n \),
       dan perhitungan perkalian skalarnya serupa.
     \end{answer}
  \item 
    Tinjau isomorfisme \( \map{\text{Rep}_B}{\polyspace_3}{\Re^4} \).
    \begin{exparts}
      \partsitem Vektor-vektor dalam ruang riil ortogonal jika dan hanya jika
        hasil kali titiknya nol.
        Berikan definisi ortogonalitas bagi polinomial.
      \partsitem Turunan suatu unsur \( \polyspace_3 \) berada dalam
        \( \polyspace_3 \).
        Berikan definisi turunan suatu vektor dalam \( \Re^4 \).
    \end{exparts}
    \begin{answer}
       \begin{exparts}
        \partsitem Menarik kembali definisi dari \( \Re^4 \) ke
          \( \polyspace_3 \) memberikan bahwa
          \( a_0+a_1x+a_2x^2+a_3x^3 \) ortogonal terhadap
          \( b_0+b_1x+b_2x^2+b_3x^3 \) jika dan hanya jika
          \( a_0b_0+a_1b_1+a_2b_2+a_3b_3=0 \).
        \partsitem Definisi yang alami adalah sebagai berikut.
          \begin{equation*}
            D(\colvec{a_0 \\ a_1 \\ a_2 \\ a_3})=
              \colvec{a_1 \\ 2a_2 \\ 3a_3 \\ 0}
          \end{equation*}
      \end{exparts}   
    \end{answer}
  \recommended \item
    Apakah setiap korespondensi antara basis, ketika diperluas ke ruang-ruangnya,
    memberikan isomorfisme?
    Artinya, andaikan \( V \) merupakan ruang vektor dengan basis
    \( B=\sequence{\vec{\beta}_1,\ldots,\vec{\beta}_n} \) dan
    \( \map{f}{B}{W} \)
    merupakan korespondensi sedemikian sehingga
    \( D=\sequence{f(\vec{\beta}_1),\ldots,f(\vec{\beta}_n)} \)
    merupakan basis bagi~$W$.
    Haruskah $\map{\hat{f}}{V}{W}$ yang membawa
    $\vec{v}=c_1\vec{\beta}_1+\cdots+c_n\vec{\beta}_n$
    ke $
    \hat{f}(\vec{v})=c_1f(\vec{\beta}_1)+\cdots+c_nf(\vec{\beta}_n)$
    merupakan isomorfisme?
    \begin{answer}
      Ya.

      Pertama, \( \hat{f} \) terdefinisi dengan baik karena setiap unsur
      \( V \) memiliki tepat satu representasi sebagai kombinasi linear
      unsur-unsur \( B \).

      Kedua, kita harus menunjukkan bahwa $\hat{f}$ injektif dan surjektif.
      Pemetaan itu injektif karena setiap unsur \( W \) hanya memiliki satu
      representasi sebagai kombinasi linear unsur-unsur \( D \), sebab $D$
      merupakan basis.
      Selanjutnya, \( \hat{f} \) surjektif karena setiap unsur \( W \)
      memiliki sekurang-kurangnya satu representasi sebagai kombinasi linear
      unsur-unsur \( D \).

      Terakhir, pemeriksaan bahwa struktur dipertahankan bersifat rutin.
      Sebagai contoh, berikut perhitungan untuk penjumlahan.
      \begin{multline*}
        \hat{f}(\,(c_1\vec{\beta}_1+\dots+c_n\vec{\beta}_n)+
                  (d_1\vec{\beta}_1+\dots+d_n\vec{\beta}_n)\,)    \\
       \begin{aligned}
        &=\hat{f}(\,(c_1+d_1)\vec{\beta}_1+\dots+(c_n+d_n)\vec{\beta}_n\,) \\
        &=(c_1+d_1)f(\vec{\beta}_1)+\dots+(c_n+d_n)f(\vec{\beta}_n) \\
        &=c_1f(\vec{\beta}_1)+\dots+c_nf(\vec{\beta}_n)
          +d_1f(\vec{\beta}_1)+\dots+d_nf(\vec{\beta}_n) \\
        &=\hat{f}(c_1\vec{\beta}_1+\dots+c_n\vec{\beta}_n)
          +\hat{f}(d_1\vec{\beta}_1+\dots+d_n\vec{\beta}_n)
      \end{aligned}
      \end{multline*}
      (Kesamaan kedua merupakan definisi $\hat{f}$.)
      Perkalian skalar diperiksa dengan cara serupa.
    \end{answer}
  \item 
    \textit{(Memerlukan subbagian opsional tentang Mengombinasikan Subruang.)}
    Andaikan \( V=V_1\directsum V_2 \) dan \( V \) isomorfik dengan ruang
    \( U \) di bawah pemetaan \( f \).
    Tunjukkan bahwa \( U=f(V_1)\directsum f(V_2) \).
    \begin{answer}
      Karena \( V_1\intersection V_2=\set{\zero_V} \) dan \( f \) injektif,
      kita memiliki \( f(V_1)\intersection f(V_2)=\set{\zero_U} \).
      Sebagai penutup, hitung dimensinya:
      \( \dim(U)=\dim(V)=\dim(V_1)+\dim(V_2)=\dim(f(V_1))+\dim(f(V_2)) \),
      sesuai kebutuhan.
    \end{answer}
  \item Tunjukkan bahwa berikut bukan fungsi yang terdefinisi dengan baik dari
    bilangan rasional ke bilangan bulat: pasangkan setiap pecahan dengan nilai
    pembilangnya.
    \begin{answer}
      Bilangan rasional memiliki banyak representasi, misalnya
      \( 1/2=3/6 \), dan pembilangnya dapat berbeda-beda di antara
      representasi tersebut.
    \end{answer}
\index{isomorfisme|)}
\end{exercises}
